%% sol-ch07.tex — Solutions for Chapter 7: Finite State Transducers

\chapter{Finite State Transducers --- Solutions}

\begin{enumerate}[{7.}1.]

%% 7.1
\item \begin{enumerate}[(a)]
  \item \textbf{Prove $\OBAR(t,yx)=\OBAR(t,y)\cdot\OBAR(\DBAR(t,y),x)$} for any Mealy machine.

  \begin{proof}
  By induction on $|y|=k$.

  \emph{Base ($k=0$):} $\OBAR(t,\lambda x)=\OBAR(t,x)=\lambda\cdot\OBAR(t,x)=\OBAR(t,\lambda)\cdot\OBAR(\DBAR(t,\lambda),x)$.

  \emph{Step:}  Let $y=\pmb{a}y'$ with $|y'|=k$.  Then:
  \begin{align*}
    \OBAR(t,\pmb{a}y'x)
    &= \omega(t,\pmb{a})\cdot\OBAR(\delta(t,\pmb{a}),y'x) \\
    &= \omega(t,\pmb{a})\cdot\OBAR(\delta(t,\pmb{a}),y')\cdot
       \OBAR(\DBAR(\delta(t,\pmb{a}),y'),x) \\
    &= \OBAR(t,\pmb{a}y')\cdot\OBAR(\DBAR(t,\pmb{a}y'),x) \\
    &= \OBAR(t,y)\cdot\OBAR(\DBAR(t,y),x). \qed
  \end{align*}
  \end{proof}

  \item \textbf{Prove $\DBAR(s,yx)=\DBAR(\DBAR(s,y),x)$.}

  \begin{proof}
  By induction on $|y|$.

  \emph{Base:} $\DBAR(s,\lambda x)=\DBAR(s,x)=\DBAR(\DBAR(s,\lambda),x)$.

  \emph{Step:}  Let $y=\pmb{a}y'$ with $|y'|=k$.  Then
  \[
    \DBAR(s,\pmb{a}y'x)=
    \DBAR(\delta(s,\pmb{a}),y'x)=
    \DBAR(\DBAR(\delta(s,\pmb{a}),y'),x)=
    \DBAR(\DBAR(s,\pmb{a}y'),x).
  \]
  Hence $\DBAR(s,yx)=\DBAR(\DBAR(s,y),x)$. \qed
  \end{proof}
\end{enumerate}

%% 7.2
\item Let $\mu:S_A\to S_B$ be a homomorphism between Mealy machines.

\begin{enumerate}[(a)]
  \item \textbf{Prove $\mu(\DBAR_A(s,x))=\DBAR_B(\mu(s),x)$.}

  \begin{proof}
  By induction on $|x|$.

  \emph{Base:} $\mu(\DBAR_A(s,\lambda))=\mu(s)=\DBAR_B(\mu(s),\lambda)$.

  \emph{Step:}  Let $x=\pmb{a}y$ with $|y|=k$.  Then
  \[
  \begin{aligned}
  \mu(\DBAR_A(s,\pmb{a}y))
  &=\mu(\DBAR_A(\delta_A(s,\pmb{a}),y))\\
  &=\DBAR_B(\mu(\delta_A(s,\pmb{a})),y)\\
  &=\DBAR_B(\delta_B(\mu(s),\pmb{a}),y)\\
  &=\DBAR_B(\mu(s),\pmb{a}y).
  \end{aligned}
  \]
  Thus $\mu(\DBAR_A(s,x))=\DBAR_B(\mu(s),x)$. \qed
  \end{proof}

  \item \textbf{Prove $\OBAR_A(s,x)=\OBAR_B(\mu(s),x)$.}

  \begin{proof}
  By induction on $|x|$.

  \emph{Base:} $\OBAR_A(s,\lambda)=\lambda=\OBAR_B(\mu(s),\lambda)$.

  \emph{Step:}  Let $x=\pmb{a}y$ with $|y|=k$.  Then
  \[
  \begin{aligned}
  \OBAR_A(s,\pmb{a}y)
  &=\omega_A(s,\pmb{a})\cdot\OBAR_A(\delta_A(s,\pmb{a}),y)\\
  &=\omega_B(\mu(s),\pmb{a})\cdot
    \OBAR_B(\mu(\delta_A(s,\pmb{a})),y)\\
  &=\omega_B(\mu(s),\pmb{a})\cdot
    \OBAR_B(\delta_B(\mu(s),\pmb{a}),y)\\
  &=\OBAR_B(\mu(s),\pmb{a}y).
  \end{aligned}
  \]
  Thus $\OBAR_A(s,x)=\OBAR_B(\mu(s),x)$. \qed
  \end{proof}
\end{enumerate}

%% 7.3
\item \textbf{Prove Corollary~7.3.}

If there is a homomorphism $\mu:A\to B$ between connected Mealy machines, then
$A$ and $B$ are equivalent.

\begin{proof}
By Exercise~7.2(b), $\OBAR_A(s_0,x)=\OBAR_B(\mu(s_0),x)$ for all $x\in\Sigma^*$.  Since $\mu(s_{0_A})=s_{0_B}$ (by definition of homomorphism), $\OBAR_A(s_{0_A},x)=\OBAR_B(s_{0_B},x)$, i.e., $f_A(x)=f_B(x)$ for all $x$.  Hence $A\equiv B$. \qed
\end{proof}

%% 7.4
\item \textbf{Prove Corollary~7.4: $M$ is minimal iff $M$ is reduced and connected.}

\begin{proof}
($\Rightarrow$) If $M$ is minimal, then $M$ must be connected by
Corollary~7.2.  Also, $M$ must be reduced by Corollary~7.3.  Hence $M$ is both
reduced and connected.

($\Leftarrow$) Assume that $M$ is reduced and connected.  Let $A$ be any minimal
FST equivalent to $M$.  By Corollaries~7.2 and~7.3, $A$ is also reduced and
connected.  Therefore Theorem~7.5 applies, so $A\cong M$.  Isomorphic machines
have the same number of states, hence $\|S_A\|=\|S_M\|$.  Since $A$ was chosen
minimal, no equivalent machine has fewer states than $A$, and therefore no
equivalent machine has fewer states than $M$.  Thus $M$ is minimal. \qed
\end{proof}

%% 7.5
\item \textbf{Pumping lemma for FTD functions.}

\begin{enumerate}[(a)]
  \item \textbf{Statement:}  Let $f:\Sigma^*\to\Gamma^*$ be FTD, implemented by a Mealy machine $M$ with $n$ states.  Then for any $x\in\Sigma^*$ with $|x|\ge n$, there exist strings $u,v,w$ with $x=uvw$, $|uv|\le n$, $|v|\ge 1$, such that for all $k\ge 0$:
  \[
    f(uv^kw) = \OBAR(s_0,u)\cdot\OBAR(\DBAR(s_0,u),v)^k\cdot\OBAR(\DBAR(s_0,uv),w).
  \]
  That is, the output of $uv^kw$ consists of the output while processing $u$, followed by $k$ repetitions of the output while processing $v$ (each time starting from state $\DBAR(s_0,u)$), followed by the output while processing $w$ from state $\DBAR(s_0,uv)=\DBAR(s_0,u)$.

  \item \textbf{Proof:}  Since $|x|\ge n$ and $M$ has $n$ states, by the pigeonhole principle the sequence of states $s_0,\DBAR(s_0,u),\DBAR(s_0,uv),\ldots$ must repeat a state within the first $n+1$ steps.  Let $u,v,w$ be such that $uvw=x$, $v$ causes a cycle: $\DBAR(s_0,u)=\DBAR(s_0,uv)$.  Then for any $k\ge 0$, $\DBAR(s_0,uv^k)=\DBAR(s_0,u)$, so the machine processes $w$ from the same state regardless of $k$.  Thus $f(uv^kw)=\OBAR(s_0,u)\cdot\OBAR(\DBAR(s_0,u),v)^k\cdot\OBAR(\DBAR(s_0,uv),w)$.  \qed
\end{enumerate}

%% 7.6
\item \textbf{Prove Lemma~7.3 parts a--e.}

\begin{enumerate}[(a)]
  \item \textbf{Part (a):} $E_{m+1,M}$ is a refinement of $E_{mM}$.

  \begin{proof}
  If $s E_{m+1,M} t$, then by definition the states $s$ and $t$ agree on all
  input strings of length at most $m+1$.  In particular, they agree on all
  strings of length at most $m$.  Hence $s E_{mM} t$. \qed
  \end{proof}

  \item \textbf{Part (b):} $E_M$ is a refinement of $E_{mM}$.

  \begin{proof}
  If $s E_M t$, then $\OBAR(s,x)=\OBAR(t,x)$ for every $x\in\Sigma^*$.  In
  particular, the two states agree on all strings of length at most $m$, so
  $s E_{mM} t$.  Therefore $E_M\subseteq E_{mM}$. \qed
  \end{proof}

  \item \textbf{Part (c):} If $E_{mM}=E_{m+1,M}$, then
  $E_{m+k,M}=E_{mM}$ for every $k\in\NN$.

  \begin{proof}
  We use induction on $k$.  The case $k=0$ is trivial, and the case $k=1$ is the
  hypothesis.  Assume $E_{m+k,M}=E_{mM}$ for some $k\ge 1$.  By part~(a),
  $E_{m+k+1,M}\subseteq E_{m+k,M}=E_{mM}$, so only the reverse inclusion needs
  proof.

  Let $s E_{mM} t$.  Since $E_{mM}=E_{m+1,M}$, we also have $s E_{m+1,M} t$.
  Therefore, for every $\pmb{a}\in\Sigma$, the successor states
  $\delta(s,\pmb{a})$ and $\delta(t,\pmb{a})$ agree on all strings of length at
  most $m$, so
  \[
    \delta(s,\pmb{a}) E_{mM} \delta(t,\pmb{a}).
  \]
  By the induction hypothesis this implies
  \[
    \delta(s,\pmb{a}) E_{m+k,M} \delta(t,\pmb{a})
  \]
  for every $\pmb{a}$.  Hence $s$ and $t$ agree on all strings of length at most
  $m+k+1$, so $s E_{m+k+1,M} t$.  Thus $E_{mM}\subseteq E_{m+k+1,M}$, and
  therefore $E_{m+k+1,M}=E_{mM}$. \qed
  \end{proof}

  \item \textbf{Part (d):} There exists $m\le \|S\|$ such that
  $E_{mM}=E_{m+1,M}$.

  \begin{proof}
  By part~(a), the relations
  \[
    E_{0M},E_{1M},E_{2M},\ldots
  \]
  form a descending chain of refinements.  If $E_{i+1,M}\neq E_{iM}$, then the
  containment is proper, so some $E_{iM}$-class is split into at least two
  $E_{i+1,M}$-classes; hence the two adjacent relations are distinct, and the
  number of equivalence classes increases by at least one.  Since there are only
  $\|S\|$ states, there can be at most $\|S\|$ equivalence classes.  Therefore
  the chain cannot refine properly more than $\|S\|-1$ times, and for some
  $m\le \|S\|$ we must have $E_{mM}=E_{m+1,M}$. \qed
  \end{proof}

  \item \textbf{Part (e):} If $E_{mM}=E_{m+1,M}$, then $E_{mM}=E_M$.

  \begin{proof}
  By part~(b), $E_M\subseteq E_{mM}$.  For the reverse inclusion, let
  $s E_{mM} t$ and let $x\in\Sigma^*$ be arbitrary, say $|x|=k$.  By part~(c),
  \[
    E_{m+k,M}=E_{mM},
  \]
  so $s E_{m+k,M} t$.  By definition of $E_{m+k,M}$, the states $s$ and $t$
  agree on every input string of length at most $m+k$, and in particular on the
  string $x$.  Thus $\OBAR(s,x)=\OBAR(t,x)$.  Since $x$ was arbitrary, $s E_M t$.
  Hence $E_{mM}\subseteq E_M$, and therefore $E_{mM}=E_M$. \qed
  \end{proof}
\end{enumerate}

%% 7.7
\item \textbf{Prove parts of Lemma~7.4 for FST $M$.}

\begin{enumerate}[(a)]
  \item \textbf{$E_{0M}$ has one class consisting of all of $S$:}

  $E_{0M}$ is defined by $s E_{0M} t$ always (no conditions).  Hence all states are in one class: $[s]_{E_{0M}}=S$ for all $s$.

  \item \textbf{$s E_{1M} t$ iff $\forall\pmb{a}\in\Sigma,\omega(s,\pmb{a})=\omega(t,\pmb{a})$:}

  $E_{1M}$ refines $E_{0M}$ by requiring that $s$ and $t$ agree on single-letter outputs.  $s E_{1M} t$ iff $s E_{0M} t$ (trivially) and $\forall\pmb{a}\in\Sigma,\omega(s,\pmb{a})=\omega(t,\pmb{a})$.

  \item \textbf{$s E_{i+1,M} t$ iff $s E_{iM} t$ and $\forall\pmb{a}\in\Sigma,\delta(s,\pmb{a}) E_{iM} \delta(t,\pmb{a})$:}

  \begin{proof}
  By definition, $E_{i+1,M}$ refines $E_{iM}$ by additionally requiring that the successors are $E_{iM}$-equivalent.  This is the standard iterative refinement.

  $(\Rightarrow)$ If $s E_{i+1,M} t$, then $s$ and $t$ agree on all strings of length $\le i+1$.  In particular, they agree on strings of length $\le i$ (so $s E_{iM} t$), and for each $\pmb{a}$, $\delta(s,\pmb{a})$ and $\delta(t,\pmb{a})$ agree on all strings of length $\le i$ (so $\delta(s,\pmb{a}) E_{iM} \delta(t,\pmb{a})$).

  $(\Leftarrow)$ If $s E_{iM} t$ and $\delta(s,\pmb{a}) E_{iM} \delta(t,\pmb{a})$ for all $\pmb{a}$, then for any $w=\pmb{a}w'$ with $|w|\le i+1$: $\OBAR(s,\pmb{a}w')=\omega(s,\pmb{a})\cdot\OBAR(\delta(s,\pmb{a}),w')$.  Since $s E_{1M} t$ (implied by $s E_{iM} t$, $i\ge 1$): $\omega(s,\pmb{a})=\omega(t,\pmb{a})$.  And $\delta(s,\pmb{a}) E_{iM} \delta(t,\pmb{a})$ gives $\OBAR(\delta(s,\pmb{a}),w')=\OBAR(\delta(t,\pmb{a}),w')$ for $|w'|\le i$.  So $\OBAR(s,w)=\OBAR(t,w)$. \qed
  \end{proof}
\end{enumerate}

%% 7.8
\item \textbf{Prove Theorem~7.6 for Moore machines.}

For a Moore machine $A=\langle\Sigma,\Gamma,S,s_0,\delta,\omega\rangle$:

\begin{proof}
We prove that
\[
  (\forall x\in\Sigma^*)(\forall y\in\Sigma^*)(\forall t\in S) 
  \OBAR(t,yx)=\OBAR(t,y)\cdot\OBAR(\DBAR(t,y),x)
\]
by induction on $|y|$.

\emph{Base ($|y|=0$):} If $y=\lambda$, then
\[
  \OBAR(t,\lambda x)=\OBAR(t,x)=\lambda\cdot\OBAR(t,x)
  =\OBAR(t,\lambda)\cdot\OBAR(\DBAR(t,\lambda),x),
\]
since $\OBAR(t,\lambda)=\lambda$ and $\DBAR(t,\lambda)=t$.

\emph{Step:} Assume the identity holds for a word $y$.  Let $\pmb{a}\in\Sigma$.
Then, by Definition~7.17,
\[
  \OBAR(t,\pmb{a}yx)
  =\omega(\delta(t,\pmb{a}))\cdot\OBAR(\delta(t,\pmb{a}),yx).
\]
Apply the inductive hypothesis with start state $\delta(t,\pmb{a})$:
\[
  \OBAR(t,\pmb{a}yx)
  =\omega(\delta(t,\pmb{a}))\cdot\OBAR(\delta(t,\pmb{a}),y)\cdot
    \OBAR(\DBAR(\delta(t,\pmb{a}),y),x).
\]
Again by Definition~7.17, the first two factors are exactly
$\OBAR(t,\pmb{a}y)$, and by Exercise~7.1(b),
$\DBAR(\delta(t,\pmb{a}),y)=\DBAR(t,\pmb{a}y)$.  Hence
\[
  \OBAR(t,\pmb{a}yx)=\OBAR(t,\pmb{a}y)\cdot\OBAR(\DBAR(t,\pmb{a}y),x).
\]
This proves the result. \qed
\end{proof}

%% 7.9
\item \textbf{Prove Theorem~7.7.}

\begin{proof}
Let
\[
  A=\langle\Sigma,\Gamma,S,s_0,\delta,\omega\rangle
\]
be a Moore machine, and let
\[
  M=\langle\Sigma,\Gamma,S,s_0,\delta,\omega'\rangle
\]
be the corresponding Mealy machine from Definition~7.19, where
$\omega'(s,\pmb{a})=\omega(\delta(s,\pmb{a}))$.

We prove by induction on $|x|$ that
\[
  (\forall s\in S)(\forall x\in\Sigma^*) 
  \OBAR_M(s,x)=\OBAR_A(s,x).
\]

\emph{Base:} If $x=\lambda$, then both sides equal $\lambda$.

\emph{Step:} Let $x=\pmb{a}y$.  Then
\[
  \OBAR_M(s,\pmb{a}y)
  =\omega'(s,\pmb{a})\cdot\OBAR_M(\delta(s,\pmb{a}),y)
  =\omega(\delta(s,\pmb{a}))\cdot\OBAR_M(\delta(s,\pmb{a}),y).
\]
By the inductive hypothesis applied to the state $\delta(s,\pmb{a})$,
\[
  \OBAR_M(s,\pmb{a}y)
  =\omega(\delta(s,\pmb{a}))\cdot\OBAR_A(\delta(s,\pmb{a}),y)
  =\OBAR_A(s,\pmb{a}y),
\]
by Definition~7.17.  Therefore $\OBAR_M(s,x)=\OBAR_A(s,x)$ for all $s$ and
$x$.  In particular,
$f_M(x)=\OBAR_M(s_0,x)=\OBAR_A(s_0,x)=f_A(x)$, so $M\equiv A$. \qed
\end{proof}

%% 7.10
\item \textbf{Prove Theorem~7.8.}

\begin{proof}
Let
\[
  M=\langle\Sigma,\Gamma,S,s_0,\delta,\omega\rangle
\]
be a Mealy machine, and let
\[
  A=\langle\Sigma,\Gamma,S\times\Gamma,\langle s_0,\alpha\rangle,\delta',\omega'\rangle
\]
be the corresponding Moore machine from Definition~7.20.

We prove the stronger statement
\[
  (\forall s\in S)(\forall \pmb{b}\in\Gamma)(\forall x\in\Sigma^*) 
  \OBAR_A(\langle s,\pmb{b}\rangle,x)=\OBAR_M(s,x).
\]

\emph{Base:} If $x=\lambda$, then both sides are $\lambda$.

\emph{Step:} Let $x=\pmb{a}y$.  By Definition~7.17 for Moore machines,
\[
  \OBAR_A(\langle s,\pmb{b}\rangle,\pmb{a}y)
  =\omega'(\delta'(\langle s,\pmb{b}\rangle,\pmb{a}))\cdot
    \OBAR_A(\delta'(\langle s,\pmb{b}\rangle,\pmb{a}),y).
\]
By Definition~7.20,
\[
  \delta'(\langle s,\pmb{b}\rangle,\pmb{a})
  =\langle\delta(s,\pmb{a}),\omega(s,\pmb{a})\rangle
\]
and
\[
  \omega'(\langle\delta(s,\pmb{a}),\omega(s,\pmb{a})\rangle)=\omega(s,\pmb{a}).
\]
Hence
\[
  \OBAR_A(\langle s,\pmb{b}\rangle,\pmb{a}y)
  =\omega(s,\pmb{a})\cdot
    \OBAR_A(\langle\delta(s,\pmb{a}),\omega(s,\pmb{a})\rangle,y).
\]
Apply the inductive hypothesis:
\[
  \OBAR_A(\langle s,\pmb{b}\rangle,\pmb{a}y)
  =\omega(s,\pmb{a})\cdot\OBAR_M(\delta(s,\pmb{a}),y)
  =\OBAR_M(s,\pmb{a}y).
\]
Therefore $\OBAR_A(\langle s,\pmb{b}\rangle,x)=\OBAR_M(s,x)$ for all states and
inputs.  In particular,
\[
  f_A(x)=\OBAR_A(\langle s_0,\alpha\rangle,x)=\OBAR_M(s_0,x)=f_M(x),
\]
so $A\equiv M$. \qed
\end{proof}

%% 7.11
\item \textbf{Use Lemma~7.6 to find $E_C$ in Example~7.10.}

For the Moore machine $C$ of Example~7.10,
\[
  \omega(r_0)=\omega(r_2)=\pmb{0},\qquad
  \omega(r_1)=\omega(r_3)=\pmb{1}.
\]
Hence
\[
  E_{0C}=\bigl\{\{r_0,r_2\},\{r_1,r_3\}\bigr\}.
\]
Now apply Lemma~7.6(b).  The states $r_0$ and $r_2$ are separated at the next
stage because
\[
  \delta(r_0,\pmb{a})=r_0,\qquad \delta(r_2,\pmb{a})=r_1,
\]
and $r_0$ and $r_1$ are in different $E_{0C}$-classes.  Likewise, $r_1$ and
$r_3$ are separated because
\[
  \delta(r_1,\pmb{a})=r_0,\qquad \delta(r_3,\pmb{a})=r_1,
\]
again landing in different $E_{0C}$-classes.  Therefore
\[
  E_{1C}=\bigl\{\{r_0\},\{r_1\},\{r_2\},\{r_3\}\bigr\}.
\]
Since $E_{1C}$ is already the identity relation, the process has stabilized and
\[
  E_C=E_{1C}.
\]

%% 7.12
\item \textbf{Homomorphism from machine $M$ (Example~7.11) to machine $B$ (Example~7.2).}

Let the states of $M$ be $\{r_0,r_1,r_2,r_3\}$, as in Example~7.11, and let
the states of $B$ be $\{s_0,s_1\}$.  Define
\[
  \mu(r_0)=\mu(r_1)=s_0,\qquad \mu(r_2)=\mu(r_3)=s_1.
\]
This is the natural choice because in Example~7.11 the states $r_0,r_1$ print
only $\pmb{0}$, while $r_2,r_3$ print only $\pmb{1}$, exactly as $s_0$ and
$s_1$ do in Example~7.2.

We check the homomorphism conditions.
\begin{enumerate}[(i)]
  \item $\mu(r_0)=s_0$, so start states are preserved.
  \item For $\pmb{a}$, Example~7.11 gives
  \[
    \delta_M(r_0,\pmb{a})=\delta_M(r_1,\pmb{a})=r_0,\qquad
    \delta_M(r_2,\pmb{a})=\delta_M(r_3,\pmb{a})=r_1.
  \]
  Therefore
  \[
    \mu(\delta_M(r_i,\pmb{a}))=s_0=\delta_B(\mu(r_i),\pmb{a})
  \]
  for every $i$, since $\delta_B(s_0,\pmb{a})=\delta_B(s_1,\pmb{a})=s_0$ in
  Example~7.2.

  For $\pmb{b}$, Example~7.11 gives
  \[
    \delta_M(r_0,\pmb{b})=\delta_M(r_1,\pmb{b})=r_2,\qquad
    \delta_M(r_2,\pmb{b})=\delta_M(r_3,\pmb{b})=r_3.
  \]
  Hence
  \[
    \mu(\delta_M(r_i,\pmb{b}))=s_1=\delta_B(\mu(r_i),\pmb{b})
  \]
  for every $i$, since $\delta_B(s_0,\pmb{b})=\delta_B(s_1,\pmb{b})=s_1$.
  \item The output table of Example~7.11 gives
  \[
    \omega_M(r_0,\pmb{a})=\omega_M(r_0,\pmb{b})
    =\omega_M(r_1,\pmb{a})=\omega_M(r_1,\pmb{b})=\pmb{0},
  \]
  and
  \[
    \omega_M(r_2,\pmb{a})=\omega_M(r_2,\pmb{b})
    =\omega_M(r_3,\pmb{a})=\omega_M(r_3,\pmb{b})=\pmb{1}.
  \]
  These agree exactly with the outputs of $B$ at $s_0$ and $s_1$.
\end{enumerate}
Thus $\mu$ is a homomorphism from $M$ onto $B$.

%% 7.13
\item \textbf{Prove $\OBAR(t,\pmb{a})=\omega(t,\pmb{a})$ for any FST.}

\begin{proof}
By definition of $\OBAR$ for a single letter:
$\OBAR(t,\pmb{a})=\omega(t,\pmb{a})\cdot\OBAR(\delta(t,\pmb{a}),\lambda)
=\omega(t,\pmb{a})\cdot\lambda=\omega(t,\pmb{a})$.
(Since $\OBAR(\delta(t,\pmb{a}),\lambda)=\lambda$ and $\lambda$ is the identity element for concatenation.) \qed
\end{proof}

%% 7.14
\item \textbf{Modified vending machine with coin return.}

Add to the machine:
\begin{itemize}
  \item New input $\pmb{r}$ (coin return).
  \item New output $\pmb{a}$ (release coins).
  \item From each state $s_k$ (representing $k$ cents in holding), the transition on $\pmb{r}$ goes to $s_0$ (reset to zero) and outputs $\pmb{a}$.
  \item The output $\pmb{a}$ could encode how many coins to return (one per 5 cents, etc.), but if $\pmb{a}$ is a single symbol meaning ``release all,'' then it's a single output.
\end{itemize}

Formally: add $\delta(s_k,\pmb{r})=s_0$ and $\omega(s_k,\pmb{r})=\pmb{a}$ for each state $s_k$ in the vending machine.

%% 7.15
\item \textbf{$\delta_{E_M}$ is well defined.}

\begin{proof}
$\delta_{E_M}([s]_{E_M},\pmb{a})=[\delta(s,\pmb{a})]_{E_M}$.  Suppose $[s]_{E_M}=[s']_{E_M}$, i.e., $s E_M s'$.  By Lemma~7.3c, $\delta(s,\pmb{a}) E_M \delta(s',\pmb{a})$.  Hence $[\delta(s,\pmb{a})]_{E_M}=[\delta(s',\pmb{a})]_{E_M}$, so $\delta_{E_M}$ gives the same result regardless of the representative chosen.
Also, for every $[s]_{E_M}\in S_{E_M}$ and every $\pmb{a}\in\Sigma$, $\delta$ is
defined on $(s,\pmb{a})$, so $[\delta(s,\pmb{a})]_{E_M}$ is an element of
$S_{E_M}$.  Thus $\delta_{E_M}$ is defined on all of
$S_{E_M}\times\Sigma$. \qed
\end{proof}

%% 7.16
\item \textbf{$\omega_{E_M}$ is well defined.}

\begin{proof}
$\omega_{E_M}([s]_{E_M},\pmb{a})=\omega(s,\pmb{a})$.  Suppose $s E_M s'$.  By the definition of $E_M$ (states are equivalent iff they produce the same output on all inputs), in particular on the single letter $\pmb{a}$: $\omega(s,\pmb{a})=\omega(s',\pmb{a})$.  Hence $\omega_{E_M}$ is independent of the choice of representative.
Also, for every $[s]_{E_M}\in S_{E_M}$ and every $\pmb{a}\in\Sigma$, $\omega$ is
defined on $(s,\pmb{a})$, so $\omega_{E_M}$ is defined on all of
$S_{E_M}\times\Sigma$. \qed
\end{proof}

%% 7.17
\item \textbf{Reduced is not sufficient for minimality.}

\textbf{Example.}  A reduced machine may have unreachable states.  Let $M$ have states $\{s_0,s_1\}$ where $s_1$ is unreachable from $s_0$.  If $s_0$ and $s_1$ have distinct output behaviors, $M$ is reduced (no two states are equivalent).  But $s_1$ is useless; removing it gives an equivalent machine with fewer states.  Hence $M$ is reduced but not minimal.

%% 7.18
\item \textbf{$\mu$ in Theorem~7.5 is well defined.}

\begin{proof}
$\mu$ is defined by $\mu(s)=\DBAR_2(s_{0_2},x)$ where $x$ is any word with $\DBAR_1(s_{0_1},x)=s$ (such $x$ exists because $M_1$ is connected).  Well-definedness requires: if $\DBAR_1(s_{0_1},x)=\DBAR_1(s_{0_1},x')=s$, then $\DBAR_2(s_{0_2},x)=\DBAR_2(s_{0_2},x')$.

Since $M_1$ and $M_2$ are equivalent, $\OBAR_1(s_{0_1},wx)=\OBAR_2(s_{0_2},wx)$ and $\OBAR_1(s_{0_1},wx')=\OBAR_2(s_{0_2},wx')$ for all $w$.  In particular, $\OBAR_1(s_{0_1},xy)=\OBAR_2(s_{0_2},xy)$ implies $\OBAR_1(s,y)=\OBAR_2(\DBAR_2(s_{0_2},x),y)$ for all $y$.  Also, $\OBAR_1(s_{0_1},x'y)=\OBAR_2(s_{0_2},x'y)$ implies $\OBAR_1(s,y)=\OBAR_2(\DBAR_2(s_{0_2},x'),y)$ for all $y$.  Since $M_2$ is reduced, the state $\DBAR_2(s_{0_2},x)$ is uniquely determined by its output behavior, which equals that of $s$.  Hence $\DBAR_2(s_{0_2},x)=\DBAR_2(s_{0_2},x')$, and $\mu$ is well defined. \qed
\end{proof}

%% 7.19
\item \textbf{$\mu$ is an isomorphism.}

\begin{enumerate}[(a)]
  \item $\mu(s_{0_1})=\DBAR_2(s_{0_2},\lambda)=s_{0_2}$.
  \item $\mu(\delta_1(s,\pmb{a}))=\DBAR_2(s_{0_2},x\pmb{a})=\delta_2(\DBAR_2(s_{0_2},x),\pmb{a})=\delta_2(\mu(s),\pmb{a})$.
  \item $\omega_1(s,\pmb{a})=\OBAR_1(s_{0_1},x\pmb{a})|_{\text{last output}}=\OBAR_2(s_{0_2},x\pmb{a})|_{\text{last}}=\omega_2(\mu(s),\pmb{a})$.  (The output functions agree since the machines are equivalent and $\mu$ commutes with transitions.)
  \item \textbf{Injective:} If $\mu(s)=\mu(t)$, then $\DBAR_2(s_{0_2},x_s)=\DBAR_2(s_{0_2},x_t)$, so $s$ and $t$ have the same output behavior from $s_{0_1}$; by $M_1$ being reduced, $s=t$.
  \item \textbf{Surjective:} Since $M_2$ is connected, every $t\in S_2$ is reachable: $t=\DBAR_2(s_{0_2},y)$ for some $y$.  Let $s=\DBAR_1(s_{0_1},y)$; then $\mu(s)=t$.
\end{enumerate}

%% 7.20
\item \textbf{One-unit delay Mealy and Moore machines.}

$\OBAR(x_1x_2\cdots x_m)=\pmb{x}\cdot x_1\cdots x_{m-1}$ (output is $\pmb{x}$ followed by the input shifted by one).

\begin{enumerate}[(a)]
  \item \textbf{Mealy sextuple:}
  $M=\langle\{\pmb{a},\pmb{b}\},\{\pmb{a},\pmb{b},\pmb{x}\},\{s_a,s_b,s_0\},s_0,\delta,\omega\rangle$ where $s_0$ is the initial state, $s_a$ means ``last input was $\pmb{a}$'', $s_b$ means ``last input was $\pmb{b}$''.
  $\delta(s_0,\pmb{a})=s_a,\ \delta(s_0,\pmb{b})=s_b,\ \delta(s_a,\pmb{a})=s_a,\ \delta(s_a,\pmb{b})=s_b,\ \delta(s_b,\pmb{a})=s_a,\ \delta(s_b,\pmb{b})=s_b$.
  $\omega(s_0,\pmb{a})=\pmb{x},\ \omega(s_0,\pmb{b})=\pmb{x},\ \omega(s_a,\pmb{a})=\pmb{a},\ \omega(s_a,\pmb{b})=\pmb{a},\ \omega(s_b,\pmb{a})=\pmb{b},\ \omega(s_b,\pmb{b})=\pmb{b}$.

  \item \textbf{Mealy machine diagram:} Three states; $s_0$ outputs $\pmb{x}$ on both inputs; $s_a$ outputs $\pmb{a}$ on both inputs; $s_b$ outputs $\pmb{b}$ on both inputs.

  \item \textbf{Moore sextuple:}  A Moore machine cannot use the same three
  states, because under Definition~7.17 the symbol printed after reading the
  first input is the output attached to the \emph{destination} state.  A
  correct Moore machine is
  \[
    A=\langle\{\pmb{a},\pmb{b}\},\{\pmb{a},\pmb{b},\pmb{x}\},Q,q_0,\delta',\eta\rangle,
  \]
  where
  \[
    Q=\{q_0,q_{x,a},q_{x,b},q_{a,a},q_{a,b},q_{b,a},q_{b,b}\}.
  \]
  Here the first component records the symbol to be printed \emph{now}, and the
  second component records the most recent input symbol for the next step.  The
  output function is
  \[
    \eta(q_0)=\pmb{x},\qquad
    \eta(q_{x,a})=\eta(q_{x,b})=\pmb{x},
  \]
  \[
    \eta(q_{a,a})=\eta(q_{a,b})=\pmb{a},\qquad
    \eta(q_{b,a})=\eta(q_{b,b})=\pmb{b}.
  \]
  The transition function is
  \[
    \delta'(q_0,\pmb{a})=q_{x,a},\qquad
    \delta'(q_0,\pmb{b})=q_{x,b},
  \]
  and, for $u\in\{\pmb{x},\pmb{a},\pmb{b}\}$ and
  $v,\sigma\in\{\pmb{a},\pmb{b}\}$,
  \[
    \delta'(q_{u,v},\sigma)=q_{v,\sigma}.
  \]
  Thus the first output is $\pmb{x}$, and each later output is the previous
  input symbol.

  \item \textbf{Moore machine diagram:}  The diagram is the graph of the
  sextuple in part~(c): the start state is $q_0/\pmb{x}$; it goes to
  $q_{x,a}/\pmb{x}$ on input $\pmb{a}$ and to $q_{x,b}/\pmb{x}$ on input
  $\pmb{b}$; and every state $q_{u,v}$ has arrows
  \[
    q_{u,v}\xrightarrow{\pmb{a}}q_{v,\pmb{a}},\qquad
    q_{u,v}\xrightarrow{\pmb{b}}q_{v,\pmb{b}}.
  \]
\end{enumerate}

%% 7.21
\item \textbf{Circuit for vending machine (Example~7.1).}

\begin{enumerate}[(a)]
  \item \textbf{EOS:} No.  The machine operates continuously, and the candy bar is
  dispensed on the transition where the button is pushed.

  \item \textbf{SOS:} Yes.  As the text points out, a vending machine should reset
  to its zero-credit state after power is restored.

  \item \textbf{Input encoding:} We want codes for
  $\{\pmb{b},\pmb{d},\pmb{n},\pmb{q},\<SOS\>\}$, so three bits suffice.
  One convenient assignment is
  \[
    \pmb{b}=000,\qquad \pmb{d}=001,\qquad \pmb{n}=010,\qquad
    \pmb{q}=011,\qquad \<SOS\>=111,
  \]
  with $100,101,110$ unused.

  \item \textbf{Output encoding:} The output alphabet is
  \[
    \{\pmb{c}_0,\pmb{c}_1,\pmb{c}_2,\pmb{c}_3,\pmb{c}_4,\pmb{d'},\pmb{n'},
    \pmb{q'},\pmb{\varphi}\},
  \]
  so four bits are needed.  Let
  \[
    \pmb{c}_0=0000,\ \pmb{c}_1=0001,\ \pmb{c}_2=0010,\ \pmb{c}_3=0011,\ 
    \pmb{c}_4=0100,\ \pmb{d'}=0101,\ \pmb{n'}=0110,\ \pmb{q'}=0111,\ 
    \pmb{\varphi}=1000.
  \]

  \item \textbf{State encoding:} Write $s_i$ for the state representing
  $5i$\cent\ for $0\leq i\leq 10$.  Then 11 states require four bits; for
  example
  \[
    s_0=0000,\ s_1=0001,\ \ldots,\ s_{10}=1010,
  \]
  with $1011,1100,1101,1110,1111$ unused.

  \item \textbf{$t_2$ truth table and Boolean function:}
  \[
  \begin{array}{c|ccccc}
  & \pmb{b} & \pmb{d} & \pmb{n} & \pmb{q} & \<SOS\> \\
  \hline
  s_0   & 0&1&0&0&0\\
  s_1   & 0&1&1&1&0\\
  s_2   & 1&0&1&1&0\\
  s_3   & 1&0&0&0&0\\
  s_4   & 0&1&0&0&0\\
  s_5   & 0&1&1&1&0\\
  s_6   & 0&1&1&1&0\\
  s_7   & 0&1&1&1&0\\
  s_8   & 0&0&0&0&0\\
  s_9   & 0&0&0&0&0\\
  s_{10}& 0&1&1&1&0
  \end{array}
  \]
  and one minimized sum-of-products is
  \[
  \begin{aligned}
  t_2'={}&(\neg t_4\wedge \neg t_2\wedge t_1\wedge \neg a_3\wedge a_2)\ \vee\
          (t_3\wedge t_2\wedge \neg a_3\wedge a_2)\ \vee\\
        &(\neg t_4\wedge \neg t_3\wedge t_2\wedge \neg a_2\wedge \neg a_1)\ \vee\
          (t_4\wedge t_2\wedge \neg a_3\wedge a_1)\ \vee\\
        &(t_2\wedge \neg t_1\wedge \neg a_3\wedge a_2)\ \vee\
          (t_3\wedge \neg a_2\wedge a_1)\ \vee\
          (\neg t_4\wedge \neg t_2\wedge \neg a_2\wedge a_1).
  \end{aligned}
  \]

  \item \textbf{$w_2$ truth table and Boolean function:}
  \[
  \begin{array}{c|ccccc}
  & \pmb{b} & \pmb{d} & \pmb{n} & \pmb{q} & \<SOS\> \\
  \hline
  s_0   & 0&0&0&0&0\\
  s_1   & 0&0&0&0&0\\
  s_2   & 0&0&0&0&0\\
  s_3   & 0&0&0&0&0\\
  s_4   & 0&0&0&0&0\\
  s_5   & 0&0&0&0&0\\
  s_6   & 0&0&1&1&0\\
  s_7   & 0&0&1&1&0\\
  s_8   & 1&0&1&1&0\\
  s_9   & 1&0&1&1&0\\
  s_{10}& 0&0&1&1&0
  \end{array}
  \]
  and a minimized form is
  \[
    w_2=(t_4\wedge \neg t_2\wedge \neg a_1)\ \vee\
        (t_3\wedge t_2\wedge \neg a_3\wedge a_2)\ \vee\
        (t_4\wedge \neg a_3\wedge a_2).
  \]

  \item The complete circuit uses four D flip-flops for the state bits
  $t_1,t_2,t_3,t_4$, three input lines $a_1,a_2,a_3$, and four output networks
  $w_1,w_2,w_3,w_4$.  The full transition/output table is exactly the one read
  from Figure~7.2: for $0\leq i\leq 5$, pushing $\pmb{b}$ leaves $s_i$ fixed with
  output $\pmb{\varphi}$, while $\pmb{n},\pmb{d},\pmb{q}$ send $s_i$ to
  $s_{i+1},s_{i+2},s_{i+5}$ with output $\pmb{\varphi}$; for $6\leq i\leq 10$,
  pushing $\pmb{b}$ sends $s_i$ to $s_0$ with output $\pmb{c}_{i-6}$, while
  $\pmb{n},\pmb{d},\pmb{q}$ self-loop with outputs $\pmb{n'},\pmb{d'},\pmb{q'}$;
  and $\<SOS\>$ sends every real state to $s_0$ with output
  $\pmb{\varphi}$.  The other next-state bits and output bits are obtained from
  this same table in exactly the same way.
\end{enumerate}

%% 7.22
\item \textbf{Circuit for modified vending machine (Exercise~7.14).}

\begin{enumerate}[(a)]
  \item \textbf{EOS:} No.

  \item \textbf{Input encoding:}  We again keep a reset code, so use three bits
  for
  $\{\pmb{b},\pmb{d},\pmb{n},\pmb{q},\pmb{r},\<SOS\>\}$:
  \[
    \pmb{b}=000,\qquad \pmb{d}=001,\qquad \pmb{n}=010,\qquad
    \pmb{q}=011,\qquad \pmb{r}=100,\qquad \<SOS\>=111.
  \]

  \item \textbf{Output encoding:}  There are now ten outputs
  $\{\pmb{a},\pmb{c}_0,\pmb{c}_1,\pmb{c}_2,\pmb{c}_3,\pmb{c}_4,\pmb{d'},
  \pmb{n'},\pmb{q'},\pmb{\varphi}\}$, so four bits suffice.  One convenient code
  is
  \[
    \pmb{a}=0000,\ \pmb{c}_0=0001,\ \pmb{c}_1=0010,\ \pmb{c}_2=0011,\ 
    \pmb{c}_3=0100,\ \pmb{c}_4=0101,\ \pmb{d'}=0110,\ \pmb{n'}=0111,\ 
    \pmb{q'}=1000,\ \pmb{\varphi}=1001.
  \]

  \item \textbf{State encoding:}  Keep the state encoding from Exercise~7.21:
  $s_i$ means $5i$\cent\ for $0\leq i\leq 10$, with
  $s_0=0000,\ldots,s_{10}=1010$.

  \item \textbf{$t_3$ truth table and Boolean function:}
  \[
  \begin{array}{c|cccccc}
  & \pmb{b} & \pmb{d} & \pmb{n} & \pmb{q} & \pmb{r} & \<SOS\> \\
  \hline
  s_0   & 0&0&0&1&0&0\\
  s_1   & 0&0&0&1&0&0\\
  s_2   & 0&1&0&1&0&0\\
  s_3   & 0&1&1&0&0&0\\
  s_4   & 1&1&1&0&0&0\\
  s_5   & 1&1&1&0&0&0\\
  s_6   & 0&1&1&1&0&0\\
  s_7   & 0&1&1&1&0&0\\
  s_8   & 0&0&0&0&0&0\\
  s_9   & 0&0&0&0&0&0\\
  s_{10}& 0&0&0&0&0&0
  \end{array}
  \]
  and one minimized form is
  \[
  \begin{aligned}
  t_3'={}&(\neg t_4\wedge t_2\wedge \neg t_1\wedge \neg a_3\wedge a_1)\ \vee\
          (t_2\wedge t_1\wedge a_2\wedge \neg a_1)\ \vee\\
        &(t_3\wedge a_2\wedge \neg a_1)\ \vee\
          (\neg t_4\wedge t_2\wedge \neg a_2\wedge a_1)\ \vee\
          (t_3\wedge t_2\wedge \neg a_3\wedge a_2)\ \vee\\
        &(t_3\wedge \neg t_2\wedge \neg a_3\wedge \neg a_2)\ \vee\
          (\neg t_4\wedge \neg t_3\wedge \neg t_2\wedge \neg a_3\wedge a_2\wedge a_1).
  \end{aligned}
  \]

  \item \textbf{$w_3$ truth table and Boolean function:}
  \[
  \begin{array}{c|cccccc}
  & \pmb{b} & \pmb{d} & \pmb{n} & \pmb{q} & \pmb{r} & \<SOS\> \\
  \hline
  s_0   & 0&0&0&0&0&0\\
  s_1   & 0&0&0&0&0&0\\
  s_2   & 0&0&0&0&0&0\\
  s_3   & 0&0&0&0&0&0\\
  s_4   & 0&0&0&0&0&0\\
  s_5   & 0&0&0&0&0&0\\
  s_6   & 0&1&1&0&0&0\\
  s_7   & 0&1&1&0&0&0\\
  s_8   & 0&1&1&0&0&0\\
  s_9   & 1&1&1&0&0&0\\
  s_{10}& 1&1&1&0&0&0
  \end{array}
  \]
  and a minimized form is
  \[
  \begin{aligned}
  w_3={}&(t_4\wedge \neg a_2\wedge a_1)\ \vee\
         (t_4\wedge t_1\wedge \neg a_3\wedge \neg a_1)\ \vee\
         (t_3\wedge t_2\wedge \neg a_2\wedge a_1)\ \vee\\
       &(t_3\wedge t_2\wedge a_2\wedge \neg a_1)\ \vee\
         (t_4\wedge a_2\wedge \neg a_1)\ \vee\
         (t_4\wedge t_2\wedge \neg a_3\wedge \neg a_1).
  \end{aligned}
  \]

  \item The complete circuit is obtained from Exercise~7.21 by adding one new
  input column for $\pmb{r}$ and changing the output table accordingly.  Every
  real state goes to $s_0$ on $\pmb{r}$ and produces $\pmb{a}$; every real state
  goes to $s_0$ on $\<SOS\>$ and produces $\pmb{\varphi}$; and all
  $\pmb{b},\pmb{d},\pmb{n},\pmb{q}$ entries are unchanged from
  Exercise~7.21.  The other state and output bits are obtained from this same
  combined table.
\end{enumerate}

%% 7.23
\item \textbf{Circuit for Example~7.2.}

Let $a_1=0$ code $\pmb{a}$ and $a_1=1$ code $\pmb{b}$; let $t_1=0$ code $s_0$
and $t_1=1$ code $s_1$; and let $w_1=0,1$ code the output symbols
$\pmb{0},\pmb{1}$.  Then
\[
\begin{array}{cc|c}
t_1 & a_1 & t_1'\\
\hline
0&0&0\\
0&1&1\\
1&0&0\\
1&1&1
\end{array}
\qquad
\begin{array}{cc|c}
t_1 & a_1 & w_1\\
\hline
0&0&0\\
0&1&0\\
1&0&1\\
1&1&1
\end{array}
\]
so the minimized equations are
\[
  t_1'=a_1,\qquad w_1=t_1.
\]
Thus the circuit uses one D flip-flop holding $t_1$, the input line $a_1$
feeding the D input, and the present-state line $t_1$ feeding the output line
$w_1$.

%% 7.24
\item \textbf{Circuit for Example~7.6.}

Encode the four states of Figure~7.4 by
\[
  t_0=00,\qquad t_1=01,\qquad t_2=10,\qquad t_3=11,
\]
let $a_1=0$ code $\pmb{a}$ and $a_1=1$ code $\pmb{b}$, and let $w_1=0,1$ code
the output symbols $\pmb{0},\pmb{1}$.  To avoid confusion with the state names,
write the two state bits as $u_2u_1$.

The truth tables are
\[
\begin{array}{ccc|cc}
u_2&u_1&a_1&u_2'&u_1'\\
\hline
0&0&0&0&1\\
0&0&1&0&0\\
0&1&0&1&0\\
0&1&1&0&0\\
1&0&0&1&0\\
1&0&1&1&1\\
1&1&0&0&1\\
1&1&1&0&0
\end{array}
\qquad
\begin{array}{ccc|c}
u_2&u_1&a_1&w_1\\
\hline
0&0&0&0\\
0&0&1&0\\
0&1&0&0\\
0&1&1&0\\
1&0&0&0\\
1&0&1&1\\
1&1&0&0\\
1&1&1&0
\end{array}
\]
and the minimized equations are
\[
  u_2'=(u_2\wedge \neg u_1)\ \vee\ (\neg u_2\wedge u_1\wedge \neg a_1),
\]
\[
  u_1'=(\neg u_2\wedge \neg u_1\wedge \neg a_1)\ \vee\
        (u_2\wedge \neg u_1\wedge a_1)\ \vee\
        (u_2\wedge u_1\wedge \neg a_1),
\]
\[
  w_1=u_2\wedge \neg u_1\wedge a_1.
\]
So the circuit has two D flip-flops holding $u_2,u_1$, a combinational network
for $u_2',u_1'$, and one output gate computing $w_1$.

%% 7.25
\item \textbf{Circuit for FST $D$ (Example~7.8).}

Encode the states of Figure~7.5 by
\[
  q_0=00,\qquad q_1=01,\qquad q_2=10,
\]
leaving $11$ as a don't-care code.  Let $a_1=0$ code $\pmb{a}$ and $a_1=1$ code
$\pmb{b}$, and let $u_2u_1$ again denote the state bits.  Then
\[
\begin{array}{ccc|cc}
u_2&u_1&a_1&u_2'&u_1'\\
\hline
0&0&0&0&1\\
0&0&1&0&0\\
0&1&0&1&0\\
0&1&1&0&0\\
1&0&0&1&0\\
1&0&1&0&0
\end{array}
\qquad
\begin{array}{ccc|c}
u_2&u_1&a_1&w_1\\
\hline
0&0&0&0\\
0&0&1&0\\
0&1&0&0\\
0&1&1&0\\
1&0&0&0\\
1&0&1&1
\end{array}
\]
with minimized equations
\[
  u_2'=\neg a_1\wedge (u_1\vee u_2),\qquad
  u_1'=\neg u_2\wedge \neg u_1\wedge \neg a_1,\qquad
  w_1=u_2\wedge a_1.
\]
Thus the circuit again has two D flip-flops and one output gate.

%% 7.26
\item \textbf{Connected is not sufficient for minimality.}

\textbf{Example.}  Let $M$ have states $\{s_0,s_1\}$, $\Sigma=\{\pmb{a}\}$, $\Gamma=\{\pmb{0},\pmb{1}\}$:
$\delta(s_0,\pmb{a})=s_1,\ \delta(s_1,\pmb{a})=s_0$; $\omega(s_0,\pmb{a})=\pmb{0},\ \omega(s_1,\pmb{a})=\pmb{0}$.
Both states output $\pmb{0}$ on $\pmb{a}$ and alternate, so $s_0 E_M s_1$.  $M$ is connected (both states reachable from $s_0$) but not reduced (two equivalent states).  Hence it is not minimal.

%% 7.27
\item \textbf{Two-unit delay Mealy and Moore machines.}

$\OBAR(x_1\cdots x_m)=\pmb{xx}x_1\cdots x_{m-2}$.

\begin{enumerate}[(a)]
  \item \textbf{Mealy sextuple:}

  Let
  \[
    M=\langle\{\pmb{a},\pmb{b}\},\{\pmb{a},\pmb{b},\pmb{x}\},S,s_{xx},\delta,\omega\rangle,
  \]
  where
  \[
    S=\{s_{xx},s_{xa},s_{xb},s_{aa},s_{ab},s_{ba},s_{bb}\}.
  \]
  The state remembers the last two input symbols, using $\pmb{x}$ as a padding
  symbol before enough input has been seen.  The transition function is
  \[
    \delta(s_{uv},\sigma)=s_{v\sigma},
  \]
  where $u,v\in\{\pmb{x},\pmb{a},\pmb{b}\}$ and the only states of this form that
  occur are the seven listed above.  The output function is
  \[
    \omega(s_{uv},\sigma)=u.
  \]
  Thus the machine prints the symbol seen two steps earlier; from the start
  state $s_{xx}$ the first two outputs are $\pmb{x}$.

  \item \textbf{Mealy machine diagram:}

  The state transition diagram has the seven states above.  Every arrow into
  $s_{v\sigma}$ is labeled $\sigma/u$, where $u$ is the first component of the
  source state $s_{uv}$.  Concretely,
  \[
    s_{xx}\xrightarrow{\pmb{a}/\pmb{x}}s_{xa},\qquad
    s_{xx}\xrightarrow{\pmb{b}/\pmb{x}}s_{xb},
  \]
  \[
    s_{xa}\xrightarrow{\pmb{a}/\pmb{x}}s_{aa},\qquad
    s_{xa}\xrightarrow{\pmb{b}/\pmb{x}}s_{ab},
  \]
  \[
    s_{xb}\xrightarrow{\pmb{a}/\pmb{x}}s_{ba},\qquad
    s_{xb}\xrightarrow{\pmb{b}/\pmb{x}}s_{bb},
  \]
  and from each state $s_{uv}$ with $u,v\in\{\pmb{a},\pmb{b}\}$,
  \[
    s_{uv}\xrightarrow{\pmb{a}/u}s_{va},\qquad
    s_{uv}\xrightarrow{\pmb{b}/u}s_{vb}.
  \]

  \item \textbf{Moore sextuple:}

  A correct Moore machine must distinguish the ``first input seen'' stage from
  the ``two most recent inputs known'' stage.  One convenient choice is
  \[
    A=\langle\{\pmb{a},\pmb{b}\},\{\pmb{a},\pmb{b},\pmb{x}\},Q,q_0,\delta',\eta\rangle,
  \]
  where
  \[
    Q=\{q_0\}\cup\{q_{x,u}\mid u\in\{\pmb{a},\pmb{b}\}\}
      \cup\{q_{c,uv}\mid c\in\{\pmb{x},\pmb{a},\pmb{b}\},\
      u,v\in\{\pmb{a},\pmb{b}\}\}.
  \]
  The intended meanings are:
  $q_{x,u}$ means ``print $\pmb{x}$ and remember the one symbol $u$,'' while
  $q_{c,uv}$ means ``print $c$ and remember the last two input symbols $uv$.'' 
  The output function is
  \[
    \eta(q_0)=\pmb{x},\qquad
    \eta(q_{x,u})=\pmb{x},
  \]
  \[
    \eta(q_{c,uv})=c.
  \]
  The transition function is
  \[
    \delta'(q_0,\sigma)=q_{x,\sigma}
    \qquad(\sigma\in\{\pmb{a},\pmb{b}\}),
  \]
  \[
    \delta'(q_{x,u},\sigma)=q_{x,u\sigma}
    \qquad(u,\sigma\in\{\pmb{a},\pmb{b}\}),
  \]
  and, for $c\in\{\pmb{x},\pmb{a},\pmb{b}\}$ and
  $u,v,\sigma\in\{\pmb{a},\pmb{b}\}$,
  \[
    \delta'(q_{c,uv},\sigma)=q_{u,v\sigma}.
  \]
  This prints $\pmb{x}$ on the first two inputs and thereafter prints the symbol
  seen two steps earlier.

  \item \textbf{Moore machine diagram:}

  The diagram is the graph of the sextuple above.  The start state is
  $q_0/\pmb{x}$; from it,
  \[
    q_0\xrightarrow{\pmb{a}}q_{x,\pmb{a}},\qquad
    q_0\xrightarrow{\pmb{b}}q_{x,\pmb{b}}.
  \]
  Next,
  \[
    q_{x,u}\xrightarrow{\sigma}q_{x,u\sigma}
    \qquad(u,\sigma\in\{\pmb{a},\pmb{b}\}),
  \]
  and every state $q_{c,uv}$ has arrows
  \[
    q_{c,uv}\xrightarrow{\pmb{a}}q_{u,v\pmb{a}},\qquad
    q_{c,uv}\xrightarrow{\pmb{b}}q_{u,v\pmb{b}}.
  \]
\end{enumerate}

%% 7.28
\item \begin{enumerate}[(a)]
  \item Take $M_1=C$ from Example~7.6 and $M_2=D$ from Example~7.8.  These two
  transducers are equivalent, but $C$ has four states while $D$ has only three,
  so they cannot be isomorphic.  Thus the conclusion of Theorem~7.5 fails if
  $M_1$ is not reduced.

  \item The missing property is that the proposed map $\mu$ need not be
  one-to-one.
\end{enumerate}

%% 7.29
\item \begin{enumerate}[(a)]
  \item Let $M_2=D$ from Example~7.8.  Let $M_1$ be obtained from $D$ by adding
  one extra unreachable state $r$ with the same outgoing arrows as $q_0$:
  \[
    \delta(r,\pmb{a})=q_1,\qquad \delta(r,\pmb{b})=r,\qquad
    \omega(r,\pmb{a})=\omega(r,\pmb{b})=\pmb{0}.
  \]
  Then $M_1$ and $M_2$ are equivalent, but they are not isomorphic because
  $M_1$ has an extra unreachable state.

  \item The missing property is that $\mu$ need not even be well defined on all
  of $S_1$, since the definition of $\mu(s)$ requires a string reaching $s$ from
  the start state.
\end{enumerate}

%% 7.30
\item \begin{enumerate}[(a)]
  \item Again take $M_1=D$ from Example~7.8 and $M_2=C$ from Example~7.6.
  They are equivalent, but $C$ is not reduced and has one extra state, so the
  machines are not isomorphic.

  \item The missing property is that $\mu$ need not be onto $S_2$.
\end{enumerate}

%% 7.31
\item \begin{enumerate}[(a)]
  \item Let $M_1=D$ from Example~7.8, and let $M_2$ be obtained from $D$ by
  adjoining one extra unreachable state $r$ with
  \[
    \delta(r,\pmb{a})=q_1,\qquad \delta(r,\pmb{b})=r,\qquad
    \omega(r,\pmb{a})=\omega(r,\pmb{b})=\pmb{0}.
  \]
  Then $M_1\equiv M_2$, but the two machines are not isomorphic because $M_2$
  has an unreachable state that $M_1$ does not have.

  \item The missing property is surjectivity of $\mu$ onto $S_2$.
\end{enumerate}

%% 7.32
\item \begin{enumerate}[(a)]
  \item Let $A=C$ from Example~7.6.  The machine $C$ is connected, so $A^c=A$.
  But $C$ is not reduced, since Example~7.8 shows that it has an equivalent
  three-state quotient.  Hence $A$ is not reduced and $A^c$ is not reduced.

  \item Let $A$ be obtained from $D$ in Example~7.8 by adjoining one extra
  unreachable state $r$ satisfying
  \[
    \delta(r,\pmb{a})=q_1,\qquad \delta(r,\pmb{b})=r,\qquad
    \omega(r,\pmb{a})=\omega(r,\pmb{b})=\pmb{0}.
  \]
  Then $r$ is equivalent to $q_0$, so $A$ is not reduced.  However, the
  connected part $A^c$ throws away $r$ and is exactly $D$, which is reduced.
\end{enumerate}

%% 7.33
\item \textbf{Complete proof of Theorem~7.4.}

\begin{enumerate}[(a)]
  \item $\OBAR(t,y)=\OBAR_{E_M}([t]_{E_M},y)$:

  \begin{proof}
  By induction on $|y|$.

  \emph{Base:} If $y=\lambda$, then both sides equal $\lambda$.

  \emph{Step:} Let $y=\pmb{a}x$.  Then
  \[
    \OBAR_{E_M}([t]_{E_M},\pmb{a}x)
    =\omega_{E_M}([t]_{E_M},\pmb{a})\cdot
      \OBAR_{E_M}(\delta_{E_M}([t]_{E_M},\pmb{a}),x).
  \]
  By the definitions of $\omega_{E_M}$ and $\delta_{E_M}$, this is
  \[
    \omega(t,\pmb{a})\cdot
      \OBAR_{E_M}([\delta(t,\pmb{a})]_{E_M},x).
  \]
  Apply the inductive hypothesis to the state $\delta(t,\pmb{a})$:
  \[
    \OBAR_{E_M}([t]_{E_M},\pmb{a}x)
    =\omega(t,\pmb{a})\cdot\OBAR(\delta(t,\pmb{a}),x)
    =\OBAR(t,\pmb{a}x).
  \]
  \end{proof}

  \item $M/_{E_M}\equiv M$:

  \begin{proof}
  Part~(a) with $t=s_0$ gives
  \[
    \OBAR_{E_M}([s_0]_{E_M},y)=\OBAR(s_0,y)
  \]
  for every input string $y$.  Since $[s_0]_{E_M}$ is the start state of
  $M/_{E_M}$, the two machines define the same translation.
  \end{proof}

  \item $M/_{E_M}$ is reduced:

  \begin{proof}
  Suppose $[s]_{E_M}$ and $[t]_{E_M}$ are equivalent states of $M/_{E_M}$.  Then
  for every $y\in\Sigma^*$,
  \[
    \OBAR_{E_M}([s]_{E_M},y)=\OBAR_{E_M}([t]_{E_M},y).
  \]
  By part~(a), this implies
  \[
    \OBAR(s,y)=\OBAR(t,y)
  \]
  for every $y$.  Hence $s E_M t$, so $[s]_{E_M}=[t]_{E_M}$.  Therefore
  $M/_{E_M}$ is reduced.
  \end{proof}

  \item If $M$ connected, $M/_{E_M}$ connected:

  \begin{proof}
  Let $[t]_{E_M}$ be any state of the quotient.  Since $M$ is connected, there
  exists $x\in\Sigma^*$ with $\DBAR(s_0,x)=t$.  We claim that
  \[
    \DBAR_{E_M}([s_0]_{E_M},x)=[t]_{E_M}.
  \]
  This follows by induction on $|x|$ from the definition of $\delta_{E_M}$.
  Hence every quotient state is reachable from the quotient start state, so the
  quotient is connected.
  \end{proof}
\end{enumerate}

%% 7.34
\item \begin{enumerate}[(a)]
  \item \textbf{$f_1$ is FTD:}  Use the Mealy machine
  \[
    M=\langle\{\pmb{0},\pmb{1}\},\{\pmb{y},\pmb{n}\},
    \{\text{start},\text{read1},\text{read0}\},
    \text{start},\delta,\omega\rangle,
  \]
  where
  \[
    \delta(\text{start},\pmb{1})=\text{read1},\qquad
    \delta(\text{start},\pmb{0})=\text{read0},
  \]
  \[
    \delta(\text{read1},\pmb{0})=\delta(\text{read1},\pmb{1})=\text{read1},
  \]
  \[
    \delta(\text{read0},\pmb{0})=\delta(\text{read0},\pmb{1})=\text{read0},
  \]
  and
  \[
    \omega(\text{start},\pmb{1})=\pmb{y},\qquad
    \omega(\text{start},\pmb{0})=\pmb{n},
  \]
  \[
    \omega(\text{read1},\pmb{0})=\omega(\text{read1},\pmb{1})=\pmb{y},
  \]
  \[
    \omega(\text{read0},\pmb{0})=\omega(\text{read0},\pmb{1})=\pmb{n}.
  \]
  After the first input, the machine enters a state that remembers whether that
  first symbol was $\pmb{1}$ or $\pmb{0}$ and keeps printing $\pmb{y}$ or
  $\pmb{n}$ forever.  Hence it computes $f_1$.

  \item \textbf{$f_2$ is not FTD:}  Assume $f_2$ were FTD.  Then Exercise~7.36
  would apply: for every $n$ and every $x\in\Sigma^n$, the first $n$ letters of
  $f_2(xy)$ and $f_2(xz)$ must agree for all suffixes $y,z$.

  Take $x=\pmb{0}^n$.  Then
  \[
    f_2(\pmb{0}^n\pmb{0})=\pmb{n}^{n+1},
    \qquad
    f_2(\pmb{0}^n\pmb{1})=\pmb{y}^{n+1}.
  \]
  These two input strings agree on their first $n$ letters, but the first $n$
  letters of the outputs are completely different.  This contradicts
  Exercise~7.36.  Therefore $f_2$ is not FTD. \qed
\end{enumerate}

%% 7.35
\item \textbf{$f_3$ is not FTD.}

Assume $f_3$ were FTD, and let $M$ be a Mealy machine with $n$ states that
computes it.  Since $f_3$ depends only on the input length, we may feed $M$ the
unary inputs
\[
  \pmb{a},\pmb{a}^2,\pmb{a}^3,\ldots
\]
and examine the resulting infinite output sequence
\[
  \pmb{010010001000010^5 10^6 10^7\cdots}.
\]
Among the states reached after reading
\[
  \lambda,\pmb{a},\pmb{a}^2,\ldots,\pmb{a}^n,
\]
some state must repeat.  So there exist integers $0\le i<j\le n$ such that
\[
  \DBAR(s_0,\pmb{a}^i)=\DBAR(s_0,\pmb{a}^j).
\]
Let
\[
  u=\pmb{a}^i,\qquad v=\pmb{a}^{j-i}.
\]
Then Exercise~7.5 implies that the outputs on the strings
\[
  u,\ uv,\ uv^2,\ uv^3,\ldots
\]
are obtained by repeating one fixed output block corresponding to $v$.  Hence
the infinite output stream produced while repeatedly reading $\pmb{a}$ would be
ultimately periodic.

But the target sequence is not ultimately periodic, because the number of
consecutive $\pmb{0}$s between successive $\pmb{1}$s is
\[
  1,2,3,4,5,\ldots,
\]
which grows without bound.  Therefore $f_3$ cannot be FTD. \qed

%% 7.36
\item \textbf{Prove: for FTD $f$, the first $n$ letters of $f(xy)$ equal those of $f(xz)$.}

\begin{proof}
Let $M$ be a Mealy machine for $f$ with states $S$.  For any $x\in\Sigma^n$:
$\DBAR(s_0,x)$ is some state $s$.  Then $f(xy)=\OBAR(s_0,xy)=\OBAR(s_0,x)\cdot\OBAR(s,y)$ and $f(xz)=\OBAR(s_0,xz)=\OBAR(s_0,x)\cdot\OBAR(s,z)$.  The first $n$ letters of both are $\OBAR(s_0,x)$ (since $|x|=n$ and the first $n$ output letters come from processing $x$).  Hence they agree on the first $n$ letters. \qed
\end{proof}

%% 7.37
\item \textbf{Elevator transducer.}

\begin{enumerate}[(a)]
  \item \textbf{Input alphabet:}  Since several buttons may be active
  simultaneously, use
  \[
    \Sigma=\{0,1\}^4,
  \]
  with an input written $\langle u,d,c_1,c_2\rangle$.  Here $u$ is the hall call
  ``up from floor 1,'' $d$ is the hall call ``down from floor 2,'' and
  $c_1,c_2$ are the car buttons for floors 1 and 2.

  \item \textbf{Output alphabet:}
  \[
    \Gamma=\{\text{close},\text{open},\text{goto1},\text{goto2}\}.
  \]

  \item \textbf{Mealy sextuple:}  Let
  \[
    M=\langle \Sigma,\Gamma,S,(1,\text{closed},\emptyset),\delta,\omega\rangle,
  \]
  where
  \[
    S=\{(f,m,P)\mid f\in\{1,2\},\ m\in\{\text{open},\text{closed}\},\
    P\subseteq\{1,2\}\}.
  \]
  Here $f$ is the current floor, $m$ is the door status, and $P$ is the set of
  pending requests.  For an input $x=\langle u,d,c_1,c_2\rangle$, define
  \[
    R(x)=\{ 1\mid u=1\text{ or }c_1=1 \}\cup
         \{ 2\mid d=1\text{ or }c_2=1 \},
  \]
  and let $P^*=P\cup R(x)$.  Then:
  \begin{enumerate}[(i)]
    \item if $m=\text{open}$, then
    \[
      \delta((f,\text{open},P),x)=(f,\text{closed},P^*),\qquad
      \omega((f,\text{open},P),x)=\text{close};
    \]
    \item if $m=\text{closed}$ and $f\in P^*$, then
    \[
      \delta((f,\text{closed},P),x)=(f,\text{open},P^*\setminus\{f\}),\qquad
      \omega((f,\text{closed},P),x)=\text{open};
    \]
    \item if $m=\text{closed}$, $f=1$, and $2\in P^*$, then
    \[
      \delta((1,\text{closed},P),x)=(2,\text{closed},P^*\setminus\{2\}),\qquad
      \omega((1,\text{closed},P),x)=\text{goto2};
    \]
    \item if $m=\text{closed}$, $f=2$, and $1\in P^*$, then
    \[
      \delta((2,\text{closed},P),x)=(1,\text{closed},P^*\setminus\{1\}),\qquad
      \omega((2,\text{closed},P),x)=\text{goto1};
    \]
    \item if $m=\text{closed}$ and $P^*=\emptyset$, then
    \[
      \delta((f,\text{closed},\emptyset),x)=(f,\text{closed},\emptyset),\qquad
      \omega((f,\text{closed},\emptyset),x)=\text{close}.
    \]
  \end{enumerate}

  \item \textbf{Mealy machine diagram:}  The diagram is exactly the directed graph
  of the sextuple in part~(c): each state is a triple $(f,m,P)$, and each edge
  is labeled by the corresponding input/output pair from the five cases above.

  \item \textbf{Moore sextuple:}  A Moore version can use
  \[
    S'=\{C_{f,P},O_{f,P}\mid f\in\{1,2\},\ P\subseteq\{1,2\}\}\cup
       \{U_P,D_P\mid P\subseteq\{1,2\}\},
  \]
  where $C_{f,P}$ means ``at floor $f$, doors closed, pending set $P$,''
  $O_{f,P}$ means ``doors open,'' $U_P$ means ``travelling to floor 2,'' and
  $D_P$ means ``travelling to floor 1.''  The output function is
  \[
    \omega(C_{f,P})=\text{close},\qquad \omega(O_{f,P})=\text{open},\qquad
    \omega(U_P)=\text{goto2},\qquad \omega(D_P)=\text{goto1}.
  \]
  The transition function is the Moore analogue of part~(c):
  \[
    \delta(O_{f,P},x)=C_{f,P\cup R(x)},
  \]
  \[
    \delta(C_{1,P},x)=
    \begin{cases}
      O_{1,P^*\setminus\{1\}} & \text{if }1\in P^*,\\
      U_{P^*} & \text{if }1\notin P^*\text{ and }2\in P^*,\\
      C_{1,P^*} & \text{if }P^*=\emptyset,
    \end{cases}
  \]
  \[
    \delta(C_{2,P},x)=
    \begin{cases}
      O_{2,P^*\setminus\{2\}} & \text{if }2\in P^*,\\
      D_{P^*} & \text{if }2\notin P^*\text{ and }1\in P^*,\\
      C_{2,P^*} & \text{if }P^*=\emptyset,
    \end{cases}
  \]
  \[
    \delta(U_P,x)=C_{2,(P\cup R(x))\setminus\{2\}},\qquad
    \delta(D_P,x)=C_{1,(P\cup R(x))\setminus\{1\}}.
  \]

  \item \textbf{Moore machine diagram:}  The diagram is the directed graph of the
  Moore sextuple in part~(e).

  \item \textbf{Circuit:}  The Mealy machine in part~(c) has 16 states, so four
  state bits suffice.  The input is already given as four bits
  $\langle u,d,c_1,c_2\rangle$, and the outputs can be encoded by two bits, for
  example
  \[
    \text{close}=00,\qquad \text{open}=01,\qquad
    \text{goto1}=10,\qquad \text{goto2}=11.
  \]
  The required circuit is therefore the standard four-flip-flop sequential
  network implementing the transition and output rules from part~(c).
\end{enumerate}

%% 7.38
\item \textbf{Mealy machine for traffic signal (Example~7.16).}

Let
\[
  \Sigma=\{\langle 0,0\rangle,\langle 0,1\rangle,\langle 1,0\rangle,\langle 1,1\rangle\}
\]
on the sensor pair $\langle \alpha,\beta\rangle$.  A correct Mealy machine is
obtained from the Moore machine of Figure~7.15 by placing the output of the
destination state on each edge.  Use the state set
$\{s_0,s_1,s_2,s_3,s_4,s_5,s_6,s_7\}$ with
\[
\begin{aligned}
o(s_0)&=\langle G,R,G,R\rangle, &
o(s_1)&=\langle Y,R,Y,R\rangle,\\
o(s_2)&=\langle R,R,R,G\rangle, &
o(s_3)&=\langle G,R,Y,R\rangle,\\
o(s_4)&=\langle R,R,R,Y\rangle, &
o(s_5)&=\langle G,Y,R,R\rangle,\\
o(s_6)&=\langle G,G,R,R\rangle, &
o(s_7)&=\langle Y,Y,R,R\rangle.
\end{aligned}
\]
The transition table is
\[
\begin{array}{c|cccc}
& \langle 0,0\rangle & \langle 0,1\rangle & \langle 1,0\rangle & \langle 1,1\rangle\\
\hline
s_0 & s_0 & s_1 & s_3 & s_1\\
s_1 & s_2 & s_2 & s_2 & s_2\\
s_2 & s_4 & s_2 & s_4 & s_2\\
s_3 & s_6 & s_6 & s_6 & s_6\\
s_4 & s_0 & s_0 & s_6 & s_6\\
s_5 & s_0 & s_0 & s_0 & s_0\\
s_6 & s_5 & s_7 & s_6 & s_6\\
s_7 & s_2 & s_2 & s_2 & s_2
\end{array}
\]
and the Mealy output on any transition $(s,x)$ is
\[
  \omega(s,x)=o(\delta(s,x)).
\]

%% 7.39
\item \textbf{Traffic signal with walk signals.}

\begin{enumerate}[(a)]
  \item The new alphabets are
  \[
    \Sigma=\{0,1\}^3
  \]
  on $\langle \alpha,\beta,\gamma\rangle$, and
  \[
    \Gamma=\{R,Y,G\}^4\times\{D,W\}.
  \]

  \item \textbf{Moore machine:}  Use the states
  \[
    s_0,\ s_1^D,\ s_1^W,\ s_2^D,\ s_2^W,\ s_3,\ s_4,\ s_5,\ s_6,\ s_7,
  \]
  with outputs
  \[
  \begin{aligned}
  \omega(s_0)&=\langle G,R,G,R,D\rangle, &
  \omega(s_1^D)&=\langle Y,R,Y,R,D\rangle,\\
  \omega(s_1^W)&=\langle Y,R,Y,R,D\rangle, &
  \omega(s_2^D)&=\langle R,R,R,G,D\rangle,\\
  \omega(s_2^W)&=\langle R,R,R,G,W\rangle, &
  \omega(s_3)&=\langle G,R,Y,R,D\rangle,\\
  \omega(s_4)&=\langle R,R,R,Y,D\rangle, &
  \omega(s_5)&=\langle G,Y,R,R,D\rangle,\\
  \omega(s_6)&=\langle G,G,R,R,D\rangle, &
  \omega(s_7)&=\langle Y,Y,R,R,D\rangle.
  \end{aligned}
  \]
  The transition policy is:
  \begin{enumerate}[(i)]
    \item from $s_0$, stay in $s_0$ on $000$; go to $s_3$ on $100$; and whenever
    $\beta=1$ or $\gamma=1$, go to $s_1^W$ if $\gamma=1$ and to $s_1^D$
    otherwise;
    \item $s_1^D$ goes to $s_2^D$ on every input, and $s_1^W$ goes to $s_2^W$ on
    every input;
    \item $s_2^W$ goes to $s_2^D$ on every input, so the walk signal reverts to
    $D$ before the side-street light turns yellow;
    \item $s_2^D$ stays in $s_2^D$ while $\beta=1$ and goes to $s_4$ when
    $\beta=0$;
    \item $s_4$ goes to $s_6$ if $\alpha=1$ and to $s_0$ otherwise;
    \item $s_3$ goes to $s_6$ on every input;
    \item $s_6$ stays in $s_6$ while $\alpha=1$; if $\alpha=0$ and either
    $\beta=1$ or $\gamma=1$, it goes to $s_7$; otherwise it goes to $s_5$;
    \item $s_5$ goes to $s_0$ on every input, and $s_7$ goes to $s_2^W$ when
    $\gamma=1$ and to $s_2^D$ otherwise.
  \end{enumerate}

  \item \textbf{Mealy machine:}  Take the state set and transition graph from
  part~(b), and place on each edge the output of the destination state.  If $o(t)$ denotes the
  five-component output attached to the Moore state $t$, then
  \[
    \omega(s,x)=o(\delta(s,x)).
  \]
\end{enumerate}

%% 7.40
\item \textbf{Intersection with left-turn signals.}

Let
\[
  \Sigma=\{0,1\}^4
\]
on $\langle \alpha,\beta,\gamma,\delta\rangle$, and let
\[
  \Gamma=\{R,Y,G\}^8
\]
with coordinates ordered by signals $1,2,3,4,5,6,7,8$ in Figure~7.20.

A convenient Moore controller uses the twelve phase states
\[
  EW_G,\ EW_Y,\ NS_G,\ NS_Y,\ WB_G,\ WB_Y,\ NB_G,\ NB_Y,\ EB_G,\ EB_Y,\ SB_G,\ SB_Y,
\]
with outputs
\[
\begin{aligned}
EW_G&=\langle G,R,G,R,R,R,R,R\rangle, &
EW_Y&=\langle Y,R,Y,R,R,R,R,R\rangle,\\
NS_G&=\langle R,G,R,G,R,R,R,R\rangle, &
NS_Y&=\langle R,Y,R,Y,R,R,R,R\rangle,\\
WB_G&=\langle G,R,R,R,G,R,R,R\rangle, &
WB_Y&=\langle Y,R,R,R,Y,R,R,R\rangle,\\
NB_G&=\langle R,G,R,R,R,G,R,R\rangle, &
NB_Y&=\langle R,Y,R,R,R,Y,R,R\rangle,\\
EB_G&=\langle R,R,G,R,R,R,G,R\rangle, &
EB_Y&=\langle R,R,Y,R,R,R,Y,R\rangle,\\
SB_G&=\langle R,R,R,G,R,R,R,G\rangle, &
SB_Y&=\langle R,R,R,Y,R,R,R,Y\rangle.
\end{aligned}
\]
One explicit transition policy is:
\begin{enumerate}[(i)]
  \item start in $EW_G$; remain there while no sensor is active, and move to
  $EW_Y$ as soon as some request appears;
  \item from $EW_Y$, give priority to the two sensor-controlled lanes on the
  north-south side: go to $NB_G$ if $\beta=1$, else to $SB_G$ if $\delta=1$,
  else to $NS_G$;
  \item $NB_G$ stays in $NB_G$ while $\beta=1$, then goes to $NB_Y$, and
  $NB_Y$ goes to $SB_G$ if $\delta=1$ and otherwise to $NS_G$;
  \item $SB_G$ stays in $SB_G$ while $\delta=1$, then goes to $SB_Y$, and
  $SB_Y$ goes to $NS_G$;
  \item $NS_G$ remains in $NS_G$ while $\alpha=\gamma=0$, and moves to $NS_Y$
  once some east-west left-turn request appears;
  \item from $NS_Y$, give priority first to $\alpha$ and then to $\gamma$: go
  to $WB_G$ if $\alpha=1$, else to $EB_G$ if $\gamma=1$, else to $EW_G$;
  \item $WB_G$ stays in $WB_G$ while $\alpha=1$, then goes to $WB_Y$, and
  $WB_Y$ goes to $EB_G$ if $\gamma=1$ and otherwise to $EW_G$;
  \item $EB_G$ stays in $EB_G$ while $\gamma=1$, then goes to $EB_Y$, and
  $EB_Y$ goes to $EW_G$.
\end{enumerate}
This gives a complete Moore machine.  The corresponding Mealy machine uses the
same graph and places the output of the destination phase on each transition.

%% 7.41
\item \textbf{Base-3 adder.}

\begin{enumerate}[(a)]
  \item \textbf{Input alphabet:}
  \[
    \Sigma=\{0,1,2\}\times\{0,1,2\}.
  \]

  \item \textbf{Output alphabet:}
  \[
    \Gamma=\{0,1,2\}.
  \]

  \item \textbf{Mealy machine:}  Let
  \[
    M=\langle \Sigma,\Gamma,\{C_0,C_1\},C_0,\delta,\omega\rangle,
  \]
  where $C_0$ means ``no carry'' and $C_1$ means ``carry 1.''  For
  $\langle a,b\rangle\in\Sigma$, define
  \[
    \delta(C_i,\langle a,b\rangle)=
    \begin{cases}
      C_1 & \text{if }a+b+i\ge 3,\\
      C_0 & \text{if }a+b+i\le 2,
    \end{cases}
  \]
  and
  \[
    \omega(C_i,\langle a,b\rangle)\equiv a+b+i \pmod 3.
  \]
  This is exactly the base-3 analogue of the binary adder in Example~7.15.

  \item \textbf{Moore machine:}  A corresponding Moore machine is
  \[
    A=\langle \Sigma,\Gamma,S,q_*,\delta',\eta\rangle,
  \]
  where
  \[
    S=\{q_*\}\cup\{q_{ij}\mid i\in\{0,1\},\ j\in\{0,1,2\}\},
  \]
  the output function is
  \[
    \eta(q_*)=0,\qquad \eta(q_{ij})=j,
  \]
  and the transition function is
  \[
    \delta'(q_*,\langle a,b\rangle)=q_{c,s},
  \]
  \[
    \delta'(q_{ij},\langle a,b\rangle)=q_{c',s'},
  \]
  where
  \[
    s\equiv a+b \pmod 3,\qquad
    c=
    \begin{cases}
      1 & \text{if }a+b\ge 3,\\
      0 & \text{if }a+b\le 2,
    \end{cases}
  \]
  and
  \[
    s'\equiv a+b+i \pmod 3,\qquad
    c'=
    \begin{cases}
      1 & \text{if }a+b+i\ge 3,\\
      0 & \text{if }a+b+i\le 2.
    \end{cases}
  \]
  Thus the Moore state records both the current carry and the digit that is to
  be printed.

  \item \textbf{Circuit with \<EOS\> and \<SOS\>:}  Encode each trit by two bits,
  for example $0=00$, $1=01$, $2=10$, and encode a column
  $\langle a,b\rangle$ by the four bits of the two trits.  Reserve extra input
  codes for $\<SOS\>$ and $\<EOS\>$.  One flip-flop stores the carry bit $c$.
  The $\<SOS\>$ input forces $c=0$.  On an ordinary input column
  $\langle a,b\rangle$, the combinational network computes
  \[
    s\equiv a+b+c \pmod 3,\qquad
    c'=
    \begin{cases}
      1 & \text{if }a+b+c\ge 3,\\
      0 & \text{if }a+b+c\le 2,
    \end{cases}
  \]
  and sends the two-bit code of $s$ to the output device.  On $\<EOS\>$, the
  circuit outputs the final carry digit: print $1$ if $c=1$ and $0$ if $c=0$,
  then reset the flip-flop to $0$.  This is the required circuit implementation
  of the Mealy machine in part~(c).
\end{enumerate}

%% 7.42
\item \textbf{Base-10 adder.}

\begin{enumerate}[(a)]
  \item \textbf{Input and output alphabets:}

  \[
    \Sigma=\{0,1,\ldots,9\}\times\{0,1,\ldots,9\},\qquad
    \Gamma=\{0,1,\ldots,9\}.
  \]
  As in Example~7.15, one input symbol is a column of two digits, and the
  transducer reads the columns from right to left.

  \item \textbf{Mealy sextuple:}

  Let
  \[
    M=\langle\Sigma,\Gamma,\{N,C\},N,\delta,\omega\rangle,
  \]
  where $N$ means ``no carry'' and $C$ means ``carry.''  For
  $(i,j)\in\Sigma$, define
  \[
    \omega(N,(i,j))=(i+j)\bmod 10,\qquad
    \delta(N,(i,j))=
    \begin{cases}
      C & \text{if } i+j\ge 10,\\
      N & \text{if } i+j<10,
    \end{cases}
  \]
  and
  \[
    \omega(C,(i,j))=(i+j+1)\bmod 10,\qquad
    \delta(C,(i,j))=
    \begin{cases}
      C & \text{if } i+j+1\ge 10,\\
      N & \text{if } i+j+1<10.
    \end{cases}
  \]

  \item \textbf{Moore sextuple:}

  A convenient Moore version uses one start state $s$ together with states
  \[
    q_{c,d}\qquad(c\in\{0,1\},\ d\in\{0,1,\ldots,9\}),
  \]
  where $c$ records the carry to be used on the next column and $d$ is the most
  recently produced output digit.  Let
  \[
    A=\langle\Sigma,\Gamma,\{s\}\cup\{q_{c,d}\mid c\in\{0,1\},0\le d\le 9\},
    s,\delta',\eta\rangle.
  \]
  The output function is
  \[
    \eta(q_{c,d})=d,
  \]
  while the value of $\eta(s)$ is irrelevant because the start state's output is
  not used on empty input.  For $(i,j)\in\Sigma$, define
  \[
    \delta'(s,(i,j))=q_{c',d'},
  \]
  where
  \[
    d'=(i+j)\bmod 10,\qquad
    c'=
    \begin{cases}
      1 & \text{if } i+j\ge 10,\\
      0 & \text{if } i+j<10,
    \end{cases}
  \]
  and in general
  \[
    \delta'(q_{c,d},(i,j))=q_{c',d'},
  \]
  where
  \[
    d'=(i+j+c)\bmod 10,\qquad
    c'=
    \begin{cases}
      1 & \text{if } i+j+c\ge 10,\\
      0 & \text{if } i+j+c<10.
    \end{cases}
  \]

  \item \textbf{Circuit:}

  Using \<EOS\> and \<SOS\>, encode each decimal digit in BCD and use one extra
  carry flip-flop.  The combinational part computes
  \[
    s=i+j+c,
  \]
  outputs the BCD code for $s\bmod 10$, and feeds back the carry bit indicating
  whether $s\ge 10$.  The \<SOS\> input initializes the carry flip-flop to $0$,
  and the \<EOS\> input can be used to report whether a final overflow carry is
  present.
\end{enumerate}

%% 7.43
\item \textbf{Left-to-right addition is not FTD.}

Assume for contradiction that the left-to-right addition function $f_4$ is FTD.
Then Exercise~7.36 applies: for every $n$, whenever two inputs agree on their
first $n$ symbols, the first $n$ output symbols must also agree.

Fix any $n\ge 1$.  Consider the following two length-$(n+1)$ input strings of
column pairs:
\[
  x=\langle 0,0\rangle\langle 0,1\rangle^{n-1}\langle 1,1\rangle
  \qquad\text{and}\qquad
  y=\langle 0,0\rangle\langle 0,1\rangle^{n-1}\langle 0,1\rangle.
\]
These two inputs agree on their first $n$ symbols.  However, their first output
symbols are different when the addition is carried out correctly from right to
left:
\begin{itemize}
  \item For $x$, the last column $\langle 1,1\rangle$ generates a carry.  Each
  middle column $\langle 0,1\rangle$ propagates that carry, so a carry reaches
  the first column $\langle 0,0\rangle$.  Hence the first output bit is
  $\pmb{1}$.
  \item For $y$, the last column $\langle 0,1\rangle$ generates no carry, so no
  carry reaches the first column.  Hence the first output bit is $\pmb{0}$.
\end{itemize}
Thus the first output symbols of $f_4(x)$ and $f_4(y)$ differ, even though the
first $n$ input symbols of $x$ and $y$ are identical.  This contradicts
Exercise~7.36.  Therefore $f_4$ is not FTD.

%% 7.44
\item \textbf{Prove Theorem~7.9 for MSM.}

\begin{enumerate}[(a)]
  \item $M/_{E_M}\equiv M$:

  \begin{proof}
  As in the Mealy case, one first proves by induction on $|y|$ that for every
  state $t$ of $M$,
  \[
    \OBAR_M(t,y)=\OBAR_{E_M}([t]_{E_M},y).
  \]
  The base case is immediate since both sides are $\lambda$ on empty input.  For
  the inductive step, use the Moore-machine definition of \OBAR, together with
  the facts that
  \[
    \delta_{E_M}([s]_{E_M},\pmb{a})=[\delta(s,\pmb{a})]_{E_M}
    \quad\text{and}\quad
    \omega_{E_M}([s]_{E_M})=\omega(s).
  \]
  Taking $t=s_0$ yields
  $f_{M/_{E_M}}=f_M$, so the quotient is equivalent to $M$.
  \end{proof}

  \item $M/_{E_M}$ is reduced:

  \begin{proof}
  Suppose $[s]_{E_M}$ and $[t]_{E_M}$ are equivalent states of the quotient.
  Then
  \[
    \OBAR_{E_M}([s]_{E_M},y)=\OBAR_{E_M}([t]_{E_M},y)
  \]
  for all $y\in\Sigma^*$.  By part~(a), this implies
  $\OBAR_M(s,y)=\OBAR_M(t,y)$ for all $y$, so $s E_M t$.  Therefore
  $[s]_{E_M}=[t]_{E_M}$, and the quotient is reduced.
  \end{proof}

  \item If $M$ connected, $M/_{E_M}$ connected:

  \begin{proof}
  Let $[t]_{E_M}$ be any state of the quotient.  Since $M$ is connected, there
  exists $x\in\Sigma^*$ with $\DBAR_M(s_0,x)=t$.  By induction on $|x|$,
  \[
    \DBAR_{E_M}([s_0]_{E_M},x)=[\DBAR_M(s_0,x)]_{E_M}=[t]_{E_M}.
  \]
  Thus every quotient state is reachable, so the quotient is connected.
  \end{proof}
\end{enumerate}

%% 7.45
\item \textbf{Minimal Moore machine is reduced.}

\begin{proof}
Suppose the minimal Moore machine $M$ for $f$ is not reduced.  Then two states $s\neq t$ in $M$ have $s E_M t$.  Form $M'$ by merging $[s]=[t]$ in the quotient $M/_{E_M}$; this gives an equivalent machine with fewer states.  But $M$ was defined as minimal---contradiction.  Hence $M$ is reduced. \qed
\end{proof}

%% 7.46
\item \textbf{Theorem~7.10 for MSM.}

\begin{enumerate}[(a)]
  \item $M$ minimal iff $M$ is reduced and connected:

  \begin{proof}
  If $M$ is minimal, then it must be connected, since unreachable states can be
  deleted without changing the translation.  It must also be reduced by
  Theorem~7.9(d).  Hence minimality implies reduced and connected.

  Conversely, assume $M$ is reduced and connected.  Let $B$ be any Moore machine
  equivalent to $M$.  Replace $B$ by $B^c/_{E_{B^c}}$ if necessary.  By
  Theorem~7.9(a)--(c), this replacement is equivalent to $B$, reduced, and
  connected, and it has no more states than $B$.  So it is enough to prove that
  every reduced connected Moore machine equivalent to $M$ has at least as many
  states as $M$.

  Let $B$ now be reduced, connected, and equivalent to $M$.  Since $M$ is
  connected, for each state $s\in S_M$ choose a string $x_s\in\Sigma^*$ such
  that
  \[
    \DBAR_M(s_{0_M},x_s)=s.
  \]
  Define
  \[
    \mu(s)=\DBAR_B(s_{0_B},x_s).
  \]
  Since $M$ is connected, such a string $x_s$ exists for every $s\in S_M$; hence
  this rule assigns a value to every state of $M$.  We show that $\mu$ is well
  defined and one-to-one.

  First, suppose $x_s$ and $y_s$ both reach the same state $s$ in $M$.  For any
  $z\in\Sigma^*$, Theorem~7.6 gives
  \[
    \OBAR_M(s,z)=\OBAR_B(\DBAR_B(s_{0_B},x_s),z)
  \]
  and also
  \[
    \OBAR_M(s,z)=\OBAR_B(\DBAR_B(s_{0_B},y_s),z),
  \]
  because $M$ and $B$ are equivalent and both starts have processed the same
  prefixes $x_s$ and $y_s$.  Thus the two states of $B$ reached by $x_s$ and
  $y_s$ are equivalent.  Since $B$ is reduced, they are equal.  Hence $\mu$ is
  well defined.

  Next, suppose $\mu(s)=\mu(t)$.  Then for every $z\in\Sigma^*$,
  \[
    \OBAR_M(s,z)=\OBAR_B(\mu(s),z)=\OBAR_B(\mu(t),z)=\OBAR_M(t,z),
  \]
  so $s E_M t$.  Because $M$ is reduced, $s=t$.  Therefore $\mu$ is injective.

  Since $\mu$ is an injective map from $S_M$ into $S_B$,
  \[
    \|S_M\|\le \|S_B\|.
  \]
  Hence no equivalent Moore machine can have fewer states than $M$, and $M$ is
  minimal.
  \end{proof}

  \item $M^c/_{E_{M^c}}$ is minimal:

  \begin{proof}
  By Theorem~7.9, the quotient $M^c/_{E_{M^c}}$ is equivalent to $M^c$, reduced,
  and connected.  Since $M^c$ is equivalent to $M$, the quotient is also
  equivalent to $M$.  Part~(a) now implies that $M^c/_{E_{M^c}}$ is minimal.
  \end{proof}
\end{enumerate}

%% 7.47
\item \textbf{Prove Lemma~7.5 and Corollary~7.9 for MSMs.}

\begin{enumerate}[(a)]
  \item $(\forall s \in S_A)(\forall x \in \Sigma^*)(\mu(\DBAR_A(s,x)) =
  \DBAR_B(\mu(s),x))$.

  \begin{proof}
  By induction on $|x|$.

  \emph{Base:} If $x=\lambda$, then
  $\mu(\DBAR_A(s,\lambda))=\mu(s)=\DBAR_B(\mu(s),\lambda)$.

  \emph{Step:} Let $x=y\pmb{a}$.  Then
  \[
    \mu(\DBAR_A(s,y\pmb{a}))
    =\mu(\delta_A(\DBAR_A(s,y),\pmb{a}))
    =\delta_B(\mu(\DBAR_A(s,y)),\pmb{a}),
  \]
  by the homomorphism property.  Apply the inductive hypothesis to obtain
  \[
    \mu(\DBAR_A(s,y\pmb{a}))=\delta_B(\DBAR_B(\mu(s),y),\pmb{a})
    =\DBAR_B(\mu(s),y\pmb{a}).
  \]
  \end{proof}

  \item $(\forall s \in S_A)(\forall x \in \Sigma^*)(\OBAR_A(s,x)=
  \OBAR_B(\mu(s),x))$.

  \begin{proof}
  By induction on $|x|$.

  \emph{Base:} If $x=\lambda$, then both sides equal $\lambda$.

  \emph{Step:} Let $x=\pmb{a}y$.  By Definition~7.17,
  \[
    \OBAR_A(s,\pmb{a}y)
    =\omega_A(\delta_A(s,\pmb{a}))\cdot\OBAR_A(\delta_A(s,\pmb{a}),y).
  \]
  Since $\mu$ is a Moore homomorphism,
  \[
    \omega_A(\delta_A(s,\pmb{a}))
    =\omega_B(\mu(\delta_A(s,\pmb{a})))
    =\omega_B(\delta_B(\mu(s),\pmb{a})).
  \]
  Applying the inductive hypothesis to the state $\delta_A(s,\pmb{a})$ gives
  \[
    \OBAR_A(s,\pmb{a}y)
    =\omega_B(\delta_B(\mu(s),\pmb{a}))\cdot
      \OBAR_B(\delta_B(\mu(s),\pmb{a}),y)
    =\OBAR_B(\mu(s),\pmb{a}y).
  \]
  \end{proof}

  \item $A\equiv B$:

  \begin{proof}
  Taking $s=s_{0_A}$ in part~(b), and using $\mu(s_{0_A})=s_{0_B}$, we obtain
  \[
    \OBAR_A(s_{0_A},x)=\OBAR_B(s_{0_B},x)
  \]
  for every $x\in\Sigma^*$.  Hence $f_A=f_B$.
  \end{proof}
\end{enumerate}

%% 7.48
\item \textbf{Prove Corollary~7.10.}

\begin{proof}
Start with $E_{0M}$, which groups together exactly those states that print the
same output symbol.  By Lemma~7.6(b), compute successively
$E_{1M},E_{2M},\ldots$ from the transition function.  By Lemma~7.3(d), some
equality $E_{mM}=E_{m+1,M}$ must occur after finitely many steps, and by
Lemma~7.3(e), that stabilized relation is exactly $E_M$.  Thus there is an
algorithm for computing $E_M$. \qed
\end{proof}

%% 7.49
\item \textbf{Prove Corollary~7.11.}

\begin{proof}
Compute the connected part $M^c$ by starting with the set containing only the
start state and repeatedly adjoining every state reached by one transition from
states already found.  Then use
Corollary~7.10 to compute $E_{M^c}$.  Finally, form the quotient
$M^c/_{E_{M^c}}$.  By Theorem~7.10(b), this quotient is the minimal Moore
machine equivalent to $M$.  Hence there is an algorithm for computing the
minimal machine equivalent to $M$. \qed
\end{proof}

%% 7.50
\item \textbf{$\delta_{E_M}$ and $\omega_{E_M}$ well-defined for FST quotient.}

\begin{enumerate}[(a)]
  \item $\delta_{E_M}$ is well-defined:

  \begin{proof}
  Define
  \[
    \delta_{E_M}([s]_{E_M},\pmb{a})=[\delta(s,\pmb{a})]_{E_M}.
  \]
  If $[s]_{E_M}=[t]_{E_M}$, then $s E_M t$.  By Lemma~7.3(c),
  \[
    \delta(s,\pmb{a}) E_M \delta(t,\pmb{a}),
  \]
  so
  \[
    [\delta(s,\pmb{a})]_{E_M}=[\delta(t,\pmb{a})]_{E_M}.
  \]
  Therefore $\delta_{E_M}$ does not depend on the chosen representative.
  \end{proof}

  \item $\omega_{E_M}$ is well-defined:

  \begin{proof}
  Define
  \[
    \omega_{E_M}([s]_{E_M},\pmb{a})=\omega(s,\pmb{a}).
  \]
  If $[s]_{E_M}=[t]_{E_M}$, then $s E_M t$, so the states $s$ and $t$ agree on
  every one-letter input.  In particular,
  \[
    \omega(s,\pmb{a})=\omega(t,\pmb{a}).
  \]
  Hence $\omega_{E_M}$ is independent of the representative.
  \end{proof}
\end{enumerate}

%% 7.51
\item \textbf{$\delta_{E_M}$ and $\omega_{E_M}$ well-defined for MSM quotient.}

\begin{enumerate}[(a)]
  \item $\delta_{E_M}$ is well-defined:

  \begin{proof}
  Define
  \[
    \delta_{E_M}([s]_{E_M},\pmb{a})=[\delta(s,\pmb{a})]_{E_M}.
  \]
  Suppose $[s]_{E_M}=[t]_{E_M}$, so $s E_M t$.  Then $s$ and $t$ produce the
  same output on every continuation.  In particular, for every
  $x\in\Sigma^*$,
  \[
    \OBAR(s,\pmb{a}x)=\OBAR(t,\pmb{a}x).
  \]
  By the Moore-machine definition of \OBAR, this implies
  \[
    \omega(\delta(s,\pmb{a}))\cdot\OBAR(\delta(s,\pmb{a}),x)
    =
    \omega(\delta(t,\pmb{a}))\cdot\OBAR(\delta(t,\pmb{a}),x).
  \]
  Hence $\delta(s,\pmb{a}) E_M \delta(t,\pmb{a})$, and so
  \[
    [\delta(s,\pmb{a})]_{E_M}=[\delta(t,\pmb{a})]_{E_M}.
  \]
  Therefore $\delta_{E_M}$ does not depend on the representative chosen.
  \end{proof}

  \item $\omega_{E_M}$ is well-defined:

  \begin{proof}
  Define
  \[
    \omega_{E_M}([s]_{E_M})=\omega(s).
  \]
  If $[s]_{E_M}=[t]_{E_M}$, then $s E_M t$.  By Lemma~7.3(b), the full state
  equivalence relation $E_M$ is a refinement of $E_{0M}$, so $s E_{0M} t$.
  Lemma~7.6(a) says that $s E_{0M} t$ is exactly the statement that the two
  states have the same Moore output.  Hence
  \[
    \omega(s)=\omega(t).
  \]
  Thus $\omega_{E_M}$ is independent of the chosen representative.
  \end{proof}
\end{enumerate}

%% 7.52
\item \textbf{Isomorphism from $A$ to $B$ and $A$ connected $\Rightarrow$ $B$ connected.}

\begin{enumerate}[(a)]
  \item \textbf{Mealy machines:}  Let $\psi:S_A\to S_B$ be an isomorphism.  Every $t\in S_B$ has a preimage $s=\psi^{-1}(t)$.  Since $A$ is connected, $s=\DBAR_A(s_{0_A},x)$ for some $x$.  Then $t=\psi(s)=\psi(\DBAR_A(s_{0_A},x))=\DBAR_B(\psi(s_{0_A}),x)=\DBAR_B(s_{0_B},x)$.  So $t$ is reachable in $B$.

  \item \textbf{Moore machines:}  Let $\psi:S_A\to S_B$ be a Moore
  isomorphism.  For any state $t\in S_B$, choose $s=\psi^{-1}(t)$.  Since $A$
  is connected, there exists $x\in\Sigma^*$ with
  \[
    s=\DBAR_A(s_{0_A},x).
  \]
  Because an isomorphism commutes with the transition function,
  \[
    t=\psi(s)=\psi(\DBAR_A(s_{0_A},x))=\DBAR_B(\psi(s_{0_A}),x)=\DBAR_B(s_{0_B},x).
  \]
  Thus every state of $B$ is reachable, so $B$ is connected.
\end{enumerate}

%% 7.53
\item \textbf{Isomorphism from $A$ to $B$ and $B$ connected $\Rightarrow$ $A$ connected.}

\begin{enumerate}[(a)]
  \item \textbf{Mealy machines:}  Let $\psi:A\to B$ be an isomorphism, and let
  $s\in S_A$.  Put $t=\psi(s)\in S_B$.  Since $B$ is connected, there exists
  $x\in\Sigma^*$ with
  \[
    t=\DBAR_B(s_{0_B},x).
  \]
  Because $\psi$ is bijective, $\psi^{-1}$ exists, and because $\psi$ preserves
  transitions, $\psi^{-1}$ does as well.  Hence
  \[
    s=\psi^{-1}(t)=\psi^{-1}(\DBAR_B(s_{0_B},x))
    =\DBAR_A(\psi^{-1}(s_{0_B}),x)=\DBAR_A(s_{0_A},x).
  \]
  So every state of $A$ is reachable, and $A$ is connected.

  \item \textbf{Moore machines:}  Let $\psi:A\to B$ be a Moore isomorphism. Since
  $\psi$ is bijective and preserves the start state and all transitions, its
  inverse $\psi^{-1}:B\to A$ also preserves the start state and transitions.
  Applying the reachability argument above to $\psi^{-1}$ shows that every state
  of $A$ is reachable from $s_{0_A}$. Thus $A$ is connected.
\end{enumerate}

%% 7.54
\item \textbf{Homomorphism from $A$ to $B$ and $A$ connected does NOT imply $B$ connected.}

\begin{enumerate}[(a)]
  \item \textbf{Mealy counterexample:}  Let $A$ have state $\{s_0\}$ (connected), $\Sigma=\{\pmb{a}\}$, $\delta_A(s_0,\pmb{a})=s_0$, $\omega_A(s_0,\pmb{a})=\pmb{0}$.  Let $B$ have states $\{t_0,t_1\}$ where $t_1$ is unreachable from $t_0$; $\delta_B(t_0,\pmb{a})=t_0$, $\omega_B(t_0,\pmb{a})=\pmb{0}$, $\delta_B(t_1,\pmb{a})=t_1$, $\omega_B(t_1,\pmb{a})=\pmb{0}$.  Map $\mu(s_0)=t_0$.  Then $\mu$ is a homomorphism, $A$ is connected, but $B$ is not connected ($t_1$ unreachable).

  \item \textbf{Moore machines:} Use the same state graph and the same map
  $\mu(s_0)=t_0$, but interpret outputs as state outputs rather than transition
  outputs.  Then $\mu$ is again a homomorphism from a connected machine onto a
  machine with an unreachable state.
\end{enumerate}

%% 7.55
\item \textbf{Homomorphism from $A$ to $B$ and $B$ connected does NOT imply $A$ connected.}

\begin{enumerate}[(a)]
  \item \textbf{Counterexample:}  Let $A$ have states $\{s_0,s_1\}$ where $s_1$ is unreachable; $B$ has state $\{t_0\}$ (trivially connected).  Map $\mu(s_0)=\mu(s_1)=t_0$.  $\mu$ is a homomorphism, $B$ is connected, but $A$ is not connected.

  \item \textbf{Moore machines:} Use the same two-state domain and one-state
  codomain, but interpret outputs as state outputs.  The map sending both states
  of $A$ to the single state of $B$ is again a homomorphism, while $A$ remains
  disconnected.
\end{enumerate}

%% 7.56
\item \textbf{If $\psi:A\to B$ and $\mu:B\to A$ are isomorphisms of connected FSTs, then $\psi=\mu^{-1}$.}

\begin{proof}
$\mu\circ\psi:S_A\to S_A$ is an isomorphism (composition of isomorphisms).  For any $s\in S_A$ reachable via $x$ ($s=\DBAR_A(s_{0_A},x)$): $(\mu\circ\psi)(s)=\mu(\psi(\DBAR_A(s_{0_A},x)))=\mu(\DBAR_B(s_{0_B},x))=\DBAR_A(s_{0_A},x)=s$.  So $\mu\circ\psi=\text{id}_{S_A}$, i.e., $\mu=\psi^{-1}$. \qed
\end{proof}

%% 7.57
\item \textbf{If $A,B$ are FSTs (possibly not connected), $\psi=\mu^{-1}$ may fail.}

\textbf{Example:}  Let $A=B$ have states $\{s_0,u,v\}$, start state $s_0$,
alphabet $\{\pmb{a}\}$, and output alphabet $\{\pmb{0}\}$, with
\[
  \delta(s_0,\pmb{a})=s_0,\qquad \delta(u,\pmb{a})=u,\qquad \delta(v,\pmb{a})=v,
\]
and every transition printing $\pmb{0}$.  The states $u$ and $v$ are both
unreachable from the start state, so the machines are not connected.  Now
\[
  \psi=\text{id}_{\{s_0,u,v\}}
\]
is an isomorphism from $A$ to $B$, and so is
\[
  \mu(s_0)=s_0,\qquad \mu(u)=v,\qquad \mu(v)=u.
\]
But $\psi^{-1}=\psi\neq \mu$.  Thus the conclusion of Exercise~7.56 fails
without connectedness.

%% 7.58
\item \textbf{Three-state MSM with $E_{0A}$ having one class; can $E_{0A}\neq E_{1A}$?}

Yes, such a three-state Moore machine exists: take
\[
  A=\langle \{\pmb{a}\},\{\pmb{0}\},\{r_0,r_1,r_2\},r_0,\delta,\omega\rangle,
\]
where
\[
  \omega(r_0)=\omega(r_1)=\omega(r_2)=\pmb{0}.
\]
and
\[
  \delta(r_0,\pmb{a})=r_1,\qquad
  \delta(r_1,\pmb{a})=r_2,\qquad
  \delta(r_2,\pmb{a})=r_2.
\]
Then Lemma~7.6(a) gives
\[
  E_{0A}=\{(r_i,r_j)\mid 0\leq i,j\leq 2\},
\]
so $E_{0A}$ has exactly one equivalence class.

However, in this situation it is \emph{not} possible for $E_{0A}$ to differ from
$E_{1A}$.  By Lemma~7.6(b),
\[
  s E_{1A} t\iff s E_{0A} t\ \wedge\ (\forall \pmb{a}\in\Sigma) 
  (\delta(s,\pmb{a})E_{0A}\delta(t,\pmb{a})).
\]
If $E_{0A}$ has only one equivalence class, then every pair of successor states
is automatically $E_{0A}$-related.  Hence every pair of states is also
$E_{1A}$-related, so in fact
\[
  E_{1A}=E_{0A}.
\]

%% 7.59
\item \begin{enumerate}[(a)]
  \item Mealy $M$ not connected and $M/_{E_M}$ not connected: take
  \[
    M=\langle \{\pmb{a}\},\{\pmb{0},\pmb{1}\},\{s_0,s_1,s_2\},s_0,\delta,\omega\rangle
  \]
  where
  \[
    \delta(s_0,\pmb{a})=s_1,\qquad
    \delta(s_1,\pmb{a})=s_1,\qquad
    \delta(s_2,\pmb{a})=s_2,
  \]
  and
  \[
    \omega(s_0,\pmb{a})=\pmb{0},\qquad
    \omega(s_1,\pmb{a})=\pmb{0},\qquad
    \omega(s_2,\pmb{a})=\pmb{1}.
  \]
  The state $s_2$ is unreachable, so $M$ is not connected.  Also
  $s_2\not E_M s_0$ and $s_2\not E_M s_1$, because on the one-letter input
  $\pmb{a}$ it outputs $\pmb{1}$ while the reachable states output $\pmb{0}$.
  Hence $[s_2]$ remains a separate unreachable state in the quotient, so
  $M/_{E_M}$ is not connected.

  \item Mealy $M$ not connected but $M/_{E_M}$ connected: take
  \[
    M=\langle \{\pmb{a}\},\{\pmb{0}\},\{t_0,t_1,t_2\},t_0,\delta,\omega\rangle
  \]
  where every state loops on $\pmb{a}$ and always outputs $\pmb{0}$:
  \[
    \delta(t_j,\pmb{a})=t_j,\qquad \omega(t_j,\pmb{a})=\pmb{0}
    \qquad (j=0,1,2).
  \]
  The states $t_1$ and $t_2$ are unreachable from $t_0$, so $M$ is not
  connected.  But all three states are equivalent, since every input word
  produces the same output word $\pmb{0}^{|x|}$ no matter where we start.
  Therefore
  \[
    [t_0]=[t_1]=[t_2],
  \]
  so $M/_{E_M}$ has just one state and is therefore connected.
\end{enumerate}

%% 7.60
\item \begin{enumerate}[(a)]
  \item Moore $M$ not connected and $M/_{E_M}$ not connected: take a machine
  with states $\{r_0,r_1,r_2\}$, start state $r_0$, alphabet $\{\pmb{a}\}$,
  outputs
  \[
    \omega(r_0)=\pmb{0},\qquad \omega(r_1)=\pmb{1},\qquad \omega(r_2)=\pmb{0},
  \]
  and transitions
  \[
    \delta(r_0,\pmb{a})=r_1,\qquad
    \delta(r_1,\pmb{a})=r_1,\qquad
    \delta(r_2,\pmb{a})=r_2.
  \]
  The state $r_2$ is unreachable, so $M$ is not connected.  Also, $r_2$ is not
  equivalent to either reachable state: it outputs $\pmb{0}$ forever, while
  $r_0$ leads to $\pmb{1}$ after one step and $r_1$ already outputs $\pmb{1}$.
  Hence the quotient still has an unreachable class, so $M/_{E_M}$ is not
  connected.

  \item Moore $M$ not connected but $M/_{E_M}$ connected: take states
  $\{r_0,r_1\}$ with start state $r_0$, alphabet $\{\pmb{a}\}$, outputs
  \[
    \omega(r_0)=\omega(r_1)=\pmb{0},
  \]
  and transitions
  \[
    \delta(r_0,\pmb{a})=r_0,\qquad \delta(r_1,\pmb{a})=r_1.
  \]
  The state $r_1$ is unreachable, so $M$ is not connected.  But $r_0$ and $r_1$
  are equivalent, since both produce only $\pmb{0}$ on every input.  Therefore
  the quotient merges them into one state, and that one-state quotient is
  connected.
\end{enumerate}

%% 7.61
\item \textbf{Prove $\mu(s) E_B \mu(t) \Leftrightarrow s E_A t$ for Mealy homomorphism.}

\begin{proof}
$(\Rightarrow)$ If $\mu(s) E_B \mu(t)$: $\OBAR_B(\mu(s),x)=\OBAR_B(\mu(t),x)$ for all $x$.  By Exercise~7.2(b), $\OBAR_A(s,x)=\OBAR_B(\mu(s),x)=\OBAR_B(\mu(t),x)=\OBAR_A(t,x)$, so $s E_A t$.

$(\Leftarrow)$ If $s E_A t$: $\OBAR_A(s,x)=\OBAR_A(t,x)$ for all $x$.  So $\OBAR_B(\mu(s),x)=\OBAR_A(s,x)=\OBAR_A(t,x)=\OBAR_B(\mu(t),x)$, giving $\mu(s) E_B \mu(t)$. \qed
\end{proof}

%% 7.62
\item \textbf{Prove $\mu(s) E_B \mu(t) \Leftrightarrow s E_A t$ for Moore homomorphism.}

\begin{proof}
$(\Rightarrow)$ Suppose $\mu(s) E_B \mu(t)$.  Then
\[
  \OBAR_B(\mu(s),x)=\OBAR_B(\mu(t),x)
\]
for every $x\in\Sigma^*$.  By Exercise~7.47(b),
\[
  \OBAR_A(s,x)=\OBAR_B(\mu(s),x)
  \quad\text{and}\quad
  \OBAR_A(t,x)=\OBAR_B(\mu(t),x),
\]
so $\OBAR_A(s,x)=\OBAR_A(t,x)$ for all $x$.  Hence $s E_A t$.

$(\Leftarrow)$ Suppose $s E_A t$.  Then
\[
  \OBAR_A(s,x)=\OBAR_A(t,x)
\]
for every $x\in\Sigma^*$.  Again using Exercise~7.47(b),
\[
  \OBAR_B(\mu(s),x)=\OBAR_A(s,x)=\OBAR_A(t,x)=\OBAR_B(\mu(t),x).
\]
Therefore $\mu(s) E_B \mu(t)$. \qed
\end{proof}

%% 7.63
\item \begin{enumerate}[(a)]
  \item FST $A$ not reduced and $A^c$ not reduced:

  Let
  \[
    A=\langle\{\pmb{a}\},\{\pmb{0}\},\{s_0,s_1,s_2\},s_0,\delta,\omega\rangle,
  \]
  where
  \[
    \delta(s_0,\pmb{a})=s_1,\qquad
    \delta(s_1,\pmb{a})=s_1,\qquad
    \delta(s_2,\pmb{a})=s_2,
  \]
  and every transition prints $\pmb{0}$.  The states $s_1$ and $s_2$ are
  equivalent, so $A$ is not reduced.  Both $s_0$ and $s_1$ are reachable from
  $s_0$, and $s_0$ and $s_1$ are also equivalent, so the connected part $A^c$
  still has equivalent distinct states.  Hence $A^c$ is not reduced.

  \item FST $A$ not reduced and $A^c$ reduced:

  Let
  \[
    A=\langle\{\pmb{a}\},\{\pmb{0},\pmb{1}\},\{s_0,s_1,s_2\},s_0,\delta,\omega\rangle,
  \]
  where
  \[
    \delta(s_0,\pmb{a})=s_1,\qquad
    \delta(s_1,\pmb{a})=s_1,\qquad
    \delta(s_2,\pmb{a})=s_2,
  \]
  and
  \[
    \omega(s_0,\pmb{a})=\pmb{0},\qquad
    \omega(s_1,\pmb{a})=\omega(s_2,\pmb{a})=\pmb{1}.
  \]
  Then $s_1 E_A s_2$, so $A$ is not reduced.  But $s_2$ is unreachable, while
  the connected part $A^c=\{s_0,s_1\}$ is reduced because $s_0$ and $s_1$ have
  different output behavior.
\end{enumerate}

%% 7.64
\item \begin{enumerate}[(a)]
  \item Moore machine $A$ not reduced and $A^c$ not reduced:

  Let the states be $\{r_0,r_1,r_2\}$ with start state $r_0$, alphabet
  $\{\pmb{a}\}$, transitions
  \[
    \delta(r_0,\pmb{a})=r_1,\qquad
    \delta(r_1,\pmb{a})=r_1,\qquad
    \delta(r_2,\pmb{a})=r_2,
  \]
  and outputs
  \[
    \omega(r_0)=\omega(r_1)=\omega(r_2)=\pmb{0}.
  \]
  Then all states are equivalent, so $A$ is not reduced.  The connected part
  $A^c$ consists of $r_0$ and $r_1$, which are still equivalent, so $A^c$ is
  not reduced.

  \item Moore machine $A$ not reduced and $A^c$ reduced:

  Let the states be $\{r_0,r_1,r_2\}$ with start state $r_0$, alphabet
  $\{\pmb{a}\}$, transitions
  \[
    \delta(r_0,\pmb{a})=r_1,\qquad
    \delta(r_1,\pmb{a})=r_1,\qquad
    \delta(r_2,\pmb{a})=r_2,
  \]
  and outputs
  \[
    \omega(r_0)=\pmb{0},\qquad \omega(r_1)=\omega(r_2)=\pmb{1}.
  \]
  Then $r_1 E_A r_2$, so $A$ is not reduced.  But $r_2$ is unreachable, while
  the connected part $A^c=\{r_0,r_1\}$ is reduced because $r_0$ and $r_1$ have
  different output behavior.
\end{enumerate}

%% 7.65
\item \textbf{Isomorphism $\cong$ is an equivalence relation on Mealy machines.}

\begin{enumerate}[(a)]
  \item \textbf{Symmetry:} If $\psi:A\to B$ is an isomorphism, then $\psi$ is
  bijective and satisfies
  \[
    \psi(s_{0_A})=s_{0_B},\qquad
    \psi(\delta_A(s,\pmb{a}))=\delta_B(\psi(s),\pmb{a}),\qquad
    \omega_A(s,\pmb{a})=\omega_B(\psi(s),\pmb{a}).
  \]
  Applying $\psi^{-1}$ to these identities shows that
  \[
    \psi^{-1}(s_{0_B})=s_{0_A},\qquad
    \psi^{-1}(\delta_B(t,\pmb{a}))=\delta_A(\psi^{-1}(t),\pmb{a}),
  \]
  and
  \[
    \omega_B(t,\pmb{a})=\omega_A(\psi^{-1}(t),\pmb{a}).
  \]
  Hence $\psi^{-1}:B\to A$ is also an isomorphism.

  \item \textbf{Reflexivity:} The identity map $\mathrm{id}_{S_A}$ preserves the
  start state, transition function, and output function, so it is an
  isomorphism from $A$ to itself.

  \item \textbf{Composition:} If $\psi:A\to B$ and $\varphi:B\to C$ are
  isomorphisms, then $\varphi\circ\psi$ is bijective.  Also,
  \[
    (\varphi\circ\psi)(s_{0_A})=\varphi(s_{0_B})=s_{0_C},
  \]
  and for every state $s$ and input $\pmb{a}$,
  \[
    (\varphi\circ\psi)(\delta_A(s,\pmb{a}))
    =\varphi(\delta_B(\psi(s),\pmb{a}))
    =\delta_C((\varphi\circ\psi)(s),\pmb{a}),
  \]
  while
  \[
    \omega_A(s,\pmb{a})
    =\omega_B(\psi(s),\pmb{a})
    =\omega_C((\varphi\circ\psi)(s),\pmb{a}).
  \]
  Thus $\varphi\circ\psi$ is an isomorphism.

  \item \textbf{Equivalence:} Parts (a), (b), and (c) show that $\cong$ is
  symmetric, reflexive, and transitive.  Therefore it is an equivalence
  relation.

  \item \textbf{Homomorphism is not an equivalence relation:} Homomorphism need
  not be symmetric.  Let $A$ have two distinct states $s_0,s_1$ with the same
  transition and output behavior, and let $B$ be the one-state machine obtained
  by identifying them.  Then there is a homomorphism $A\to B$, but no inverse
  map $B\to A$ exists, so homomorphism is not symmetric.
\end{enumerate}

%% 7.66
\item \textbf{Isomorphism $\cong$ is an equivalence relation on Moore machines; homomorphism is not.}

\begin{enumerate}[(a)]
  \item \textbf{Symmetry:} If $\psi:A\to B$ is a Moore isomorphism, then
  $\psi^{-1}$ preserves the start state, transition function, and state-output
  function, so $\psi^{-1}:B\to A$ is also an isomorphism.

  \item \textbf{Reflexivity:} The identity map on the state set of a Moore
  machine is an isomorphism.

  \item \textbf{Composition:} The composition of two Moore isomorphisms again
  preserves start state, transitions, and outputs, hence is an isomorphism.

  \item \textbf{Equivalence:} Therefore Moore-machine isomorphism is reflexive,
  symmetric, and transitive.

  \item \textbf{Homomorphism is not an equivalence relation:}  It is not
  symmetric.  Let
  \[
    A=\langle\{\pmb{a}\},\{\pmb{0}\},\{r_0,r_1\},r_0,\delta_A,\omega_A\rangle
  \]
  with
  \[
    \omega_A(r_0)=\omega_A(r_1)=\pmb{0},\qquad
    \delta_A(r_0,\pmb{a})=r_1,\qquad
    \delta_A(r_1,\pmb{a})=r_1,
  \]
  and let
  \[
    B=\langle\{\pmb{a}\},\{\pmb{0}\},\{s_0\},s_0,\delta_B,\omega_B\rangle
  \]
  with
  \[
    \omega_B(s_0)=\pmb{0},\qquad \delta_B(s_0,\pmb{a})=s_0.
  \]
  The map $\mu(r_0)=\mu(r_1)=s_0$ is a homomorphism $A\to B$.  But there is no
  homomorphism $B\to A$: any such map must send $s_0$ to the start state $r_0$,
  and then
  \[
    \mu(\delta_B(s_0,\pmb{a}))=\mu(s_0)=r_0
  \]
  would have to equal
  \[
    \delta_A(r_0,\pmb{a})=r_1,
  \]
  which is impossible.  So homomorphism is not symmetric, and therefore not an
  equivalence relation.
\end{enumerate}

%% 7.67
\item \textbf{Prove homomorphism $\mu:M\to M/_{E_M}$ exists.}

\begin{proof}
Define $\mu:S\to S_{E_M}$ by $\mu(s)=[s]_{E_M}$.

We verify the homomorphism properties.
\begin{enumerate}[(i)]
  \item $\mu(s_0)=[s_0]_{E_M}$, which is the start state of the quotient.
  \item For every state $s$ and input $\pmb{a}$,
  \[
    \mu(\delta(s,\pmb{a}))
    =[\delta(s,\pmb{a})]_{E_M}
    =\delta_{E_M}([s]_{E_M},\pmb{a})
    =\delta_{E_M}(\mu(s),\pmb{a}).
  \]
  \item For every state $s$ and input $\pmb{a}$,
  \[
    \omega_{E_M}(\mu(s),\pmb{a})
    =\omega_{E_M}([s]_{E_M},\pmb{a})
    =\omega(s,\pmb{a}).
  \]
\end{enumerate}
Thus $\mu$ is a homomorphism from $M$ onto $M/_{E_M}$. \qed
\end{proof}

%% 7.68
\item \textbf{Prove homomorphism $\mu:M\to M/_{E_M}$ exists for Moore machines.}

\begin{proof}
Define $\mu(s)=[s]_{E_M}$ from the state set of $M$ onto the state set of the
quotient.

We verify the homomorphism properties.
\begin{enumerate}[(i)]
  \item $\mu(s_0)=[s_0]_{E_M}$ is the start state of the quotient.
  \item For every state $s$ and input $\pmb{a}$,
  \[
    \mu(\delta(s,\pmb{a}))
    =[\delta(s,\pmb{a})]_{E_M}
    =\delta_{E_M}([s]_{E_M},\pmb{a})
    =\delta_{E_M}(\mu(s),\pmb{a}).
  \]
  \item For every state $s$,
  \[
    \omega_{E_M}(\mu(s))
    =\omega_{E_M}([s]_{E_M})
    =\omega(s).
  \]
\end{enumerate}
Hence $\mu$ is a Moore homomorphism from $M$ onto $M/_{E_M}$. \qed
\end{proof}

%% 7.69
\item \textbf{Limit time in state $s_2$ to 3 clock cycles.}

\begin{enumerate}[(a)]
  \item Replace $s_2$ by three states $s_{2a},s_{2b},s_{2c}$, all with output
  $\langle R,R,R,G\rangle$.  The old transition $s_1\to s_2$ becomes
  $s_1\to s_{2a}$.  Then define
  \[
  \begin{array}{c|cccc}
  & \langle 0,0\rangle & \langle 0,1\rangle & \langle 1,0\rangle & \langle 1,1\rangle\\
  \hline
  s_{2a} & s_4 & s_{2b} & s_4 & s_{2b}\\
  s_{2b} & s_4 & s_{2c} & s_4 & s_{2c}\\
  s_{2c} & s_4 & s_4 & s_4 & s_4
  \end{array}
  \]
  so the configuration $\langle R,R,R,G\rangle$ can last for one, two, or three
  cycles, but never four.

  \item Replace $s_6$ by three states $s_{6a},s_{6b},s_{6c}$, all with output
  $\langle G,G,R,R\rangle$.  Every old edge entering $s_6$ now enters
  $s_{6a}$.  Then define
  \[
  \begin{array}{c|cccc}
  & \langle 0,0\rangle & \langle 0,1\rangle & \langle 1,0\rangle & \langle 1,1\rangle\\
  \hline
  s_{6a} & s_5 & s_7 & s_{6b} & s_{6b}\\
  s_{6b} & s_5 & s_7 & s_{6c} & s_{6c}\\
  s_{6c} & s_5 & s_7 & s_5 & s_7
  \end{array}
  \]
  which bounds the left-turn green in the same way.

  \item To guarantee a \emph{minimum} of two cycles and a \emph{maximum} of four
  cycles for the left-turn green, replace $s_6$ by four states
  $s_{6a},s_{6b},s_{6c},s_{6d}$, all outputting $\langle G,G,R,R\rangle$, and
  let
  \[
  \begin{array}{c|cccc}
  & \langle 0,0\rangle & \langle 0,1\rangle & \langle 1,0\rangle & \langle 1,1\rangle\\
  \hline
  s_{6a} & s_{6b} & s_{6b} & s_{6b} & s_{6b}\\
  s_{6b} & s_5 & s_7 & s_{6c} & s_{6c}\\
  s_{6c} & s_5 & s_7 & s_{6d} & s_{6d}\\
  s_{6d} & s_5 & s_7 & s_5 & s_7
  \end{array}
  \]
  so the first transition out of $s_{6a}$ is forced, the machine may remain in
  the green left-turn configuration for two, three, or four cycles, and the
  fourth cycle is the last possible one.
\end{enumerate}

%% 7.70
\item \textbf{Ensure east-west traffic gets green occasionally.}

\begin{enumerate}[(a)]
  \item Without adding states, force every excursion away from $s_0$ to return to
  $s_0$ after one service.  Concretely, remove the self-loops on $s_2$ and $s_6$,
  and redefine
  \[
    \delta(s_2,\langle i,j\rangle)=s_4\qquad\text{for all }i,j\in\{0,1\},
  \]
  \[
    \delta(s_6,\langle 0,0\rangle)=s_5,\quad
    \delta(s_6,\langle 1,0\rangle)=s_5,\quad
    \delta(s_6,\langle 0,1\rangle)=s_7,\quad
    \delta(s_6,\langle 1,1\rangle)=s_7,
  \]
  and then force the yellow states back to east-west green:
  \[
    \delta(s_4,\langle i,j\rangle)=\delta(s_5,\langle i,j\rangle)=
    \delta(s_7,\langle i,j\rangle)=s_0
  \]
  for all $i,j$.  Then the machine must revisit
  $\langle G,R,G,R\rangle$ regularly.

  \item To obtain a genuine fairness controller, pair each phase state with a
  memory set $H\subseteq\{A,B\}$, where $A$ means ``the $\alpha$ left-turn lane
  has already had green'' and $B$ means ``the $\beta$ side-street lane has
  already had green'' since the last visit to $s_0$.  Use states $(s,H)$ with
  $s\in\{s_0,\ldots,s_7\}$ and start state $(s_0,\emptyset)$.  Whenever the
  machine enters the old state $s_2$, replace $H$ by $H\cup\{B\}$; whenever it
  enters the old state $s_6$, replace $H$ by $H\cup\{A\}$; whenever it returns to
  $s_0$, reset the memory to $\emptyset$.  The yellow states now dispatch
  according to the set of sensor-driven lanes that are still owed service:
  \begin{enumerate}[(i)]
    \item from a yellow state that follows $B$-service, if $\alpha$ is active and
    $A\notin H$, go next to the $A$-branch instead of returning to $s_0$;
    otherwise return to $s_0$;
    \item from a yellow state that follows $A$-service, if $\beta$ is active and
    $B\notin H$, go next to the $B$-branch instead of returning to $s_0$;
    otherwise return to $s_0$.
  \end{enumerate}
  Thus no sensor-controlled lane can receive a second green before every other
  active sensor-controlled lane has received a first one.
\end{enumerate}

%% 7.71
\item \textbf{Prove: if two Moore machines are homomorphic, they are equivalent.}

\begin{proof}
If $\mu:A\to B$ is a homomorphism, by Exercise~7.47(c), $f_A=f_B$.  Hence $A\equiv B$. \qed
\end{proof}

%% 7.72
\item \textbf{Prove \DSIG\ is closed under any FTD function $f:\Sigma^*\to\Sigma^*$.}

\begin{proof}
Let $L\in\DSIG$, and choose a DFA
\[
  A=\langle\Sigma,Q,q_0,\delta_A,F\rangle
\]
with $L(A)=L$.  Let
\[
  M=\langle\Sigma,\Sigma,S,s_0,\delta_M,\omega_M\rangle
\]
be a Mealy machine computing the FTD function $f$.

We build an NDFA
\[
  N=\langle\Sigma,Q\times S,\{\langle q_0,s_0\rangle\},\Delta,F\times S\rangle
\]
that accepts exactly $f(L)$.  Intuitively, while $N$ reads an output word $y$,
it guesses an input word $x$ letter by letter; the $S$-component simulates the
transducer on the guessed input, and the $Q$-component simulates the DFA $A$ on
that same guessed input.

For $\langle q,s\rangle\in Q\times S$ and $\pmb{b}\in\Sigma$, define
\[
  \Delta(\langle q,s\rangle,\pmb{b})
  =\{\langle \delta_A(q,\pmb{a}),\delta_M(s,\pmb{a})\rangle
    \mid \pmb{a}\in\Sigma,\ \omega_M(s,\pmb{a})=\pmb{b}\}.
\]
So on reading an output symbol $\pmb{b}$, the NDFA chooses any input symbol
$\pmb{a}$ that would cause the transducer, in state $s$, to print $\pmb{b}$.

By induction on the input length, after reading a word $y$, the machine $N$ can
be in exactly those states
\[
  \langle \DBAR_A(q_0,x),\DBAR_M(s_0,x)\rangle
\]
for which $f(x)=y$.  Therefore
\[
  y\in L(N)
  \iff (\exists x\in\Sigma^*)(f(x)=y\ \wedge\ \DBAR_A(q_0,x)\in F)
  \iff (\exists x\in L)(f(x)=y).
\]
Hence
\[
  L(N)=f(L).
\]
Because $N$ is an NDFA, Theorem~4.1 implies that $L(N)\in\DSIG$.  Thus
$f(L)\in\DSIG$, so \DSIG\ is closed under every FTD function $f$. \qed
\end{proof}

\end{enumerate}
